% Источники: 10_Верификация_и_тестирование/DOC-10_План_верификации.md, 10_Верификация_и_тестирование/DOC-10_Верификация_сессий.md
% Генератор: docs/build_pdf.py
% Шаблон печатного представления. Содержимое подставляет docs/build_pdf.py.
\documentclass[12pt,a4paper]{article}
\usepackage{iftex}
\ifPDFTeX
  \usepackage[utf8]{inputenc}
  \usepackage{cmap}
  \usepackage[T2A]{fontenc}
  % Tempora retains Times-like text and provides Cyrillic bold/italic shapes.
  \usepackage{tempora,newtxmath}
\else
  \usepackage{fontspec}
  \IfFontExistsTF{Times New Roman}
    {\setmainfont{Times New Roman}}
    {\setmainfont{Libertinus Serif}}
  \IfFontExistsTF{Arial}{\setsansfont{Arial}}{\setsansfont{Latin Modern Sans}}
  \IfFontExistsTF{Menlo}{\setmonofont{Menlo}[Scale=MatchLowercase]}
    {\setmonofont{Latin Modern Mono}[Scale=MatchLowercase]}
\fi
\usepackage[russian]{babel}
\usepackage{geometry}
\geometry{left=25mm,right=15mm,top=20mm,bottom=20mm}
\usepackage{setspace}
\onehalfspacing
\usepackage{indentfirst}
\setlength{\parindent}{1.25cm}
\setlength{\parskip}{2pt}
\setlength{\emergencystretch}{3em}
\usepackage{graphicx,booktabs,longtable,array,tabularx,calc,etoolbox}
\newcounter{none} % Pandoc 3.9 uses this counter for unnumbered longtables.
\usepackage{enumitem}
\setlist[itemize]{leftmargin=1.25cm,itemsep=2pt,topsep=4pt}
\setlist[enumerate]{leftmargin=1.25cm,itemsep=2pt,topsep=4pt}
\newcommand{\tightlist}{\setlength{\itemsep}{0pt}\setlength{\parskip}{0pt}}
\AtBeginEnvironment{longtable}{\small\singlespacing}
\setlength{\LTpre}{8pt}
\setlength{\LTpost}{8pt}
\setlength{\tabcolsep}{5pt}
\setlength{\LTleft}{0pt}
\setlength{\LTright}{0pt}
\usepackage{xurl}
\usepackage[unicode,colorlinks=true,linkcolor=black,urlcolor=blue]{hyperref}
\usepackage{bookmark}
\hypersetup{pdftitle={План верификации программного обеспечения},pdfauthor={Анохин Г.А.},pdfsubject={Криптотренажёр mephi-12}}
\urlstyle{same}
\pagestyle{plain}
\setcounter{secnumdepth}{3}
\setcounter{tocdepth}{2}
\listfiles

\begin{document}
\hypersetup{pageanchor=false}
\begin{titlepage}
  \thispagestyle{empty}
  \centering
  {\large НИЯУ МИФИ\par}
  \vspace*{3.5cm}
  {\LARGE\bfseries План верификации программного обеспечения\par}
  \vspace{0.5cm}
  {\large DOC-10\par}
  \vspace{0.6cm}
  {\large Платформа для изучения и практики курса по криптографии\par}
  {\large Криптотренажёр mephi-12\par}
  \vspace{0.7cm}
  {\large Версия 1.1\par}
  {\large Редакция от 25.11.2026\par}
  \vspace{1cm}
  \begin{tabularx}{\textwidth}{@{}p{3.4cm}Xp{4cm}@{}}
    \toprule
    Роль & ФИО & Подпись / Дата \\
    \midrule
    Разработал & Анохин Г.А. & \\
    \addlinespace[6pt]
    Согласовано & & \\
    \addlinespace[6pt]
    Утверждено & & \\
    \bottomrule
  \end{tabularx}
  \vfill
  {\large Москва, 2026\par}
\end{titlepage}
\hypersetup{pageanchor=true}

\section*{Лист регистрации изменений}
\addcontentsline{toc}{section}{Лист регистрации изменений}
\begingroup\small\singlespacing
\noindent
\begin{tabularx}{\textwidth}{@{}p{1.2cm}p{2cm}p{2.5cm}Xp{1.8cm}@{}}
\toprule
Версия & Дата & Автор & Содержание изменения & Утвердил \\
\midrule
1.1 & 25.11.2026 & Самбьюи А.Ж. & Уточнены собственные кейсы CFG и инфраструктуры: роли, чтение на границах срока, публикация, повторы и восстановление. Исполнение не проводилось. & \\
\bottomrule
\end{tabularx}
\endgroup
\clearpage
\begingroup\small\singlespacing
\tableofcontents
\endgroup
\clearpage

\hypertarget{doc-10-план-верификации-программного-обеспечения}{}

{\def\LTcaptype{none} % do not increment counter
\begin{longtable}[]{@{}
  >{\raggedright\arraybackslash}p{(\linewidth - 2\tabcolsep) * \real{0.3200}}
  >{\raggedright\arraybackslash}p{(\linewidth - 2\tabcolsep) * \real{0.6800}}@{}}
\toprule\noalign{}
\begin{minipage}[b]{\linewidth}\raggedright
Реквизит
\end{minipage} & \begin{minipage}[b]{\linewidth}\raggedright
Значение
\end{minipage} \\
\midrule\noalign{}
\endhead
\bottomrule\noalign{}
\endlastfoot
Система & Платформа для изучения и практики курса по криптографии в МИФИ \\
Версия и дата & 1.1 от 25.11.2026 \\
Владелец документа по плану проекта & Кокорев В.Д. \\
Подготовка исходной редакции & Анохин Г.А. \\
Уточнение кейсов CFG и инфраструктуры & Самбьюи А.Ж., 25.11.2026 \\
\end{longtable}
}

\section{Задача и границы верификации}\label{ux437ux430ux434ux430ux447ux430-ux438-ux433ux440ux430ux43dux438ux446ux44b-ux432ux435ux440ux438ux444ux438ux43aux430ux446ux438ux438}

\hypertarget{задача-и-границы-верификации}{}

Верификация --- проверка того, что разрабатываемая система выполняет предъявленные к ней требования. Для каждого требования верификатор, специалист по проверке программы, находит ситуацию, в которой его выполнение можно отличить от ошибки. Такая проверочная карточка задаёт исходные условия, действие и ожидаемый результат. Верификатор прослеживает действие от начала до конца и отмечает наблюдаемый признак успеха или отказа.

Контрольное мероприятие --- работа, которую преподаватель составляет из разных видов заданий и назначает учебной группе. После публикации мероприятие появляется в расписании. Личная сессия студента --- конкретный набор выданных ему задач с собственными данными, временем начала и результатами. Верификатор устанавливает, когда студент увидит запись в расписании, когда сможет начать сессию и к какой работе будет отнесён результат. Для самостоятельной практики служит шаблон набора задач --- имя и упорядоченный список видов заданий. Модератору поручен просмотр учебного прогресса.

Вердикт --- вывод предметного проверяющего о правильности решения, диагностика --- пояснение к ошибке. Их сравнивают с учебным эталоном, то есть с заранее известным правильным результатом для выданного условия. Для развёртывания, то есть подготовки и запуска системы в рабочей среде, проверяют порядок запуска и сохранность данных. Термины предметной области собраны в {\href{../notes/methodology/Глоссарий.pdf\#глоссарий}{глоссарии проекта}}.

План охватывает учётные записи и группы, каталог шаблонов, мероприятия и расписание. Проверяются встроенные и подключаемые виды заданий, авторские варианты и планировщик --- часть системы, которая периодически обрабатывает задачи с истёкшим сроком. Проверки затрагивают также обмен запросами и ответами по протоколу HTTP, хранение данных и секретов доступа, запуск и время отклика. Состав частей программы приведён в {\href{../07\_Программные\_требования/DOC-06\_Программные\_требования.pdf\#doc-06-сводные-программные-требования}{программных требованиях DOC-06}}, наблюдаемое поведение всей платформы --- в {\href{../04\_Системные\_требования/DOC-05\_Системные\_требования.pdf\#doc-05-системные-требования}{системных требованиях DOC-05}}. Последовательности связанных операций, в которых одна часть передаёт результат другой, входят также в план интеграционного тестирования DOC-11.

Для документарной верификации карточки разбирают по текстам требований и описанию устройства системы. Прогон означал бы исполнение действий на работающей системе с записью того, что произошло. В курсовой работе готовят план и ожидаемые результаты.

\section{Материалы и распределение работы}\label{ux43cux430ux442ux435ux440ux438ux430ux43bux44b-ux438-ux440ux430ux441ux43fux440ux435ux434ux435ux43bux435ux43dux438ux435-ux440ux430ux431ux43eux442ux44b}

\hypertarget{материалы-и-распределение-работы}{}

\begin{itemize}
\tightlist
\item
  Исходные требования приведены в {\href{../02\_Техническое\_задание/Техническое\_задание.pdf\#doc-02-техническое-задание}{техническом задании}}. Основание для карточки --- действующее системное требование. Если обязанность реализуют несколько компонентов, в карточке или в таблице связей указывают также соответствующие опубликованные программные требования. Связи, для которых есть только системное определение, прослеживают непосредственно до него.
\item
  Исходные данные проверки --- заранее подготовленные учётные записи, объекты и значения, к которым возвращаются перед разбором следующего варианта. Для проверки предметного решения нужны выданное условие, независимый эталон и ответ с одной известной ошибкой. Эталон готовит специалист соответствующей предметной области. Для случая с двумя ошибками заранее отмечают, какую из них предметная проверка должна обнаружить первой. Реальные персональные данные и действующие секреты для карточек не требуются.
\item
  Кокорев ведёт план верификации и предметную часть учётных записей. Анохин отвечает за сессии и планировщик, Георгиев --- за ранцевые задачи, Портнов --- за Эль-Гамаля, Минасов --- за демонстрационные упражнения, Самбьюи --- за шаблоны наборов, модерацию и среду развёртывания. Владельцы передают редактору исходные условия и критерии по своим областям. Портнов ведёт матрицу трассировки DOC-09 --- таблицу, по которой от требования можно перейти к обеспечивающему его правилу программы и к конкретной проверке.
\item
  Для ссылки на карточку или требование служит идентификатор (ID), то есть постоянное уникальное обозначение. При разборе несоответствия участник фиксирует версию документа, ID карточки и требования, шаг, ожидаемый и наблюдаемый исходы, а также материал, по которому расхождение можно проверить. Формальная инспекция --- разбор конкретной версии документа назначенными участниками с регистрацией замечаний и решения. Замечание, принятое на таком заседании, оформляют по {\href{../06\_Процедура\_формальной\_инспекции/Процедура\_формальной\_инспекции.pdf}{процедуре DOC-04}}. Решения об утверждении принимаются по протоколу соответствующей инспекции. Условия назначения, личного старта и завершения установлены в {\href{../02\_Техническое\_задание/Техническое\_задание.pdf\#правила-доступа-и-прохождения-контрольного-мероприятия}{правилах доступа и прохождения контрольного мероприятия}}.
\end{itemize}

\section{Методы и критерии завершения}\label{ux43cux435ux442ux43eux434ux44b-ux438-ux43aux440ux438ux442ux435ux440ux438ux438-ux437ux430ux432ux435ux440ux448ux435ux43dux438ux44f}

\hypertarget{методы-и-критерии-завершения}{}

\textbf{Системный уровень.} По сценарию студента или модератора разбирают действие и видимый ответ. Роль определяет, какие операции разрешены пользователю. При отказе отдельно проверяют, изменились ли данные. Разрешение операции без подходящей роли, выдача чужого результата и выдача частично сохранённого набора нарушают критерий, даже если интерфейс вернул обычный код успеха.

\textbf{Уровень подсистем и их взаимодействия.} Для требований к типам, полям, состояниям и срокам сопоставляют работу обработчика (части программы, выполняющей операцию) с контрактом. Контракт определяет, какие данные разрешено передать и что получит отправитель при успехе или ошибке. Две параллельные операции разбирают в обоих порядках. Для восстановления выбирают точку остановки и определяют, какая запись уже зафиксирована к этому моменту. Подробные карточки и связи по этим правилам помещены в {\hyperlink{doc-10-верификация-сессий-задач-и-планировщика}{главу о сессиях и планировщике}}.

\textbf{Эксплуатационный уровень.} По схеме запуска прослеживают зависимости служб и место хранения данных. Для производительности фиксируют границы операции и условия наблюдения. Представление сценария на бумаге позволяет проверить состав измерения. Для заключения о фактическом времени отклика потребуются запись измерения и согласованный критерий.

\textbf{Условие завершения документарной проверки.} Для каждого системного требования из объёма плана есть хотя бы одна карточка, требование обозначено существующим ID, а по исходным условиям можно определить ожидаемый результат. Для опубликованных программных требований проверяется связь с соответствующим шагом. При неопределённом критерии верификатор отмечает, какие исходные данные или договорённости требуются. При противоречии между документами указывает обе версии и передаёт вопрос их владельцам. Результат по такой позиции остаётся открытым до устранения причины и повторного разбора того же случая.

Для последующей фиксации результата используют одну запись на вариант карточки. Запись содержит ID карточки, вариант, версию объекта, условия, дату и исполнителя, фактическое значение или ссылку на наблюдение, сравнение с критерием и вывод. Совпадение всех ожидаемых признаков означает, что вариант пройден. Несовпадение хотя бы одного признака означает отклонение. Когда исходные данные недоступны или наблюдения нет, вариант отмечают как непроверенный.

\section{Как читать обозначения и ответы системы}\label{ux43aux430ux43a-ux447ux438ux442ux430ux442ux44c-ux43eux431ux43eux437ux43dux430ux447ux435ux43dux438ux44f-ux438-ux43eux442ux432ux435ux442ux44b-ux441ux438ux441ux442ux435ux43cux44b}

\hypertarget{как-читать-обозначения-и-ответы-системы}{}

Префикс \texttt{TC} (test case) обозначает проверочную карточку, \texttt{SR} --- системное требование, \texttt{SWR} --- программное. В обозначении {\hyperlink{tc-auth-001}{TC-AUTH-001}} средняя часть указывает на учётные записи и доступ, число --- порядковый номер карточки. В остальных главах используются \texttt{CFG} для шаблонов и мероприятий, \texttt{SES} для сессий, \texttt{KNAP} для ранцевых задач, \texttt{ELG} для Эль-Гамаля, \texttt{BDR} для демонстрационного упражнения и \texttt{SYS} для общих свойств. Правила выдачи обозначений приведены в {\href{../notes/methodology/Методика\_идентификаторов.pdf\#методика-идентификаторов}{методике идентификаторов}}.

Клиент --- программа, которая по действию пользователя отправляет запрос. Сервер принимает его и возвращает ответ. Интерфейс взаимодействия программ называют API. Запись операции состоит из метода HTTP и пути. Метод \texttt{GET} запрашивает данные, \texttt{POST} инициирует указанное действие (например, создание или публикацию), \texttt{PUT} передаёт изменения или ответ на задачу. Имя в фигурных скобках в пути заменяют идентификатором конкретного объекта. UUID --- уникальный идентификатор объекта, который сохраняет значение при последующих обращениях.

Параметр запроса --- именованное значение, например выбранная группа. В теле запроса или ответа передают основные данные. Для структурированных данных используется JSON --- текстовая запись именованных полей и их значений. Строка записывается в кавычках, массив перечисляет значения по порядку, \texttt{true} и \texttt{false} означают «да» и «нет», \texttt{null} --- отсутствие значения. Служебный заголовок запроса может содержать сведения, которые нужны серверу до разбора тела.

Числовой статус HTTP даёт общее значение ответа. Код 200 означает успех, 201 --- создание объекта, 400 --- ошибку входных данных, 401 --- отсутствие подтверждённой личности, 403 --- недостаток прав, 404 --- объект недоступен или отсутствует, 409 --- конфликт с текущим состоянием. Если причиной отказа могут быть разные события, в карточке рядом со статусом указан машинный код из тела ответа. Полные формы запросов и ответов, в том числе машинные коды отказов, приведены в {\href{../09\_Архитектура\_и\_API/DOC-08\_Справочник\_API.pdf\#doc-08-справочник-api-платформы}{справочнике API}}.

\section{Учётные записи и учебные группы}\label{ux443ux447ux451ux442ux43dux44bux435-ux437ux430ux43fux438ux441ux438-ux438-ux443ux447ux435ux431ux43dux44bux435-ux433ux440ux443ux43fux43fux44b}

\hypertarget{учётные-записи-и-учебные-группы}{}

Учётная запись --- сохранённые сведения, по которым система узнаёт пользователя. Для входа студент использует локальные учётные данные или подтверждённую университетскую запись home.mephi.ru. Для групповых учебных операций номер группы берётся из доверенного университетского профиля. Регистрация создаёт запись. При входе сервер подтверждает личность и выдаёт токен доступа --- подписанное сообщение о пользователе и его правах. Здесь используется формат JSON Web Token (JWT). В следующих запросах клиент передаёт этот токен в заголовке \texttt{Authorization} по схеме \texttt{Bearer}, то есть предъявляет выданный токен как подтверждение права на обращение.

\subsection{TC-AUTH-001. Локальная регистрация и границы полей}\label{tc-auth-001-ux43bux43eux43aux430ux43bux44cux43dux430ux44f-ux440ux435ux433ux438ux441ux442ux440ux430ux446ux438ux44f-ux438-ux433ux440ux430ux43dux438ux446ux44b-ux43fux43eux43bux435ux439}

\hypertarget{tc-auth-001-локальная-регистрация-и-границы-полей}{}

\hypertarget{tc-auth-001}{}

\textbf{Условия.} Для локальной регистрации подготовлены свободный и занятый адреса, имя и пароль допустимого размера. Каждый граничный вариант рассматривают на отдельном входе, сохраняя остальные поля допустимыми.

\textbf{Действие.} Зарегистрировать нового студента без необязательной фамилии и повторить регистрацию по тому же адресу. Отдельно разобрать две одновременные регистрации по одному свободному адресу. Разобрать проверку длины адреса в точках 3, 4, 120, 121 символ и пароля в точках 7, 8, 255, 256 символов. Проверку длины адреса отделяют от самостоятельной проверки его синтаксической корректности. В отдельных вариантах опустить одно из обязательных полей --- адрес, имя или пароль.

\textbf{Ожидаемый результат.} Один успешный запрос создаёт одну учётную запись и сразу выдаёт токен доступа JWT в поле \texttt{accessJwt}. Обе границы каждого допустимого интервала включены. Для недопустимой длины, неверной формы адреса или пропущенного обязательного поля ожидается 400. Повторная регистрация даёт 409. При одновременной регистрации один запрос создаёт запись, второй получает 409 без токена. При отказе новый токен отсутствует, данные существующих пользователей сохраняются.

\subsection{TC-AUTH-002. Вход, пароль и токен}\label{tc-auth-002-ux432ux445ux43eux434-ux43fux430ux440ux43eux43bux44c-ux438-ux442ux43eux43aux435ux43d}

\hypertarget{tc-auth-002-вход-пароль-и-токен}{}

\hypertarget{tc-auth-002}{}

\textbf{Условия.} А и Б обозначают разных учебных пользователей. А уже зарегистрирован. Хеш пароля --- полученное из пароля проверочное значение, с которым сервер сверяет последующий ввод. BCrypt --- используемый способ получения такого значения. Подпись HMAC-SHA256 позволяет серверу по секретному ключу подтвердить подлинность токена. Даны верный и неверный пароли, действительный токен А, просроченный токен и контрольный образец токена с неподтверждаемой подписью. Для разбора хранения берут учебную запись. Значения паролей и ключа подписи в протокол не переносят.

\textbf{Действие.} Выполнить вход с верным и неверным паролями, а также по незарегистрированному адресу. Повторить после штатного перезапуска. Сопоставить представление пароля в хранилище и способ его сравнения с условиями защищённой передачи. Прочитать личную историю с действительным токеном. Повторить чтение без токена, с просроченным и контрольным токенами. По описанию выпуска JWT проверить способ подписи и единицу срока действия.

\textbf{Ожидаемый результат.} При верном пароле клиент получает JWT, при неверном пароле или незарегистрированном адресе --- 401 без токена, в том числе после штатного перезапуска. Хранится BCrypt-хеш, сервер проверяет пароль по хешу. При передаче пароль защищён от чтения посторонней стороной и отсутствует в выдаче профиля. Подтверждённые поля \texttt{userId} и \texttt{authorities} определяют пользователя и его роли. Для подписи токена используется HMAC-SHA256, срок в днях задаёт параметр окружения. Для защищённой операции действительный токен в заголовке \texttt{Authorization:\ Bearer\ \textless{}JWT\textgreater{}} даёт доступ в пределах роли его владельца. Все три варианта без действительного токена получают 401.

\subsection{TC-AUTH-003. Свой профиль и своя история}\label{tc-auth-003-ux441ux432ux43eux439-ux43fux440ux43eux444ux438ux43bux44c-ux438-ux441ux432ux43eux44f-ux438ux441ux442ux43eux440ux438ux44f}

\hypertarget{tc-auth-003-свой-профиль-и-своя-история}{}

\hypertarget{tc-auth-003}{}

\textbf{Условия.} Студенты А и Б одной группы решили задачи одинакового типа. У А сохранены верный и неверный предметные результаты, есть действующая задача и задача, завершённая по сроку без вердикта. Ответ с одним из подтверждённых вердиктов не дошёл до А. Студент В ещё не отправлял решений.

\textbf{Действие.} Прочитать личную историю под каждой из трёх учётных записей. Для завершённых наборов повторить чтение после общего срока. Сопоставить записи с UUID соответствующих задач и сессий.

\textbf{Ожидаемый результат.} Учётная запись для личной операции определяется токеном. История А содержит оба его предметных результата, в том числе после утраты ответа. Действующая задача и просрочка без вердикта в историю решений не вошли. У Б показаны только его решения, у В --- пустой список. Чтение после срока сохраняет эти данные и не расходует попытку.

\subsection{TC-AUTH-004. Подтверждение учётной записи МИФИ и номера учебной группы}\label{tc-auth-004-ux43fux43eux434ux442ux432ux435ux440ux436ux434ux435ux43dux438ux435-ux443ux447ux451ux442ux43dux43eux439-ux437ux430ux43fux438ux441ux438-ux43cux438ux444ux438-ux438-ux43dux43eux43cux435ux440ux430-ux443ux447ux435ux431ux43dux43eux439-ux433ux440ux443ux43fux43fux44b}

\hypertarget{tc-auth-004-подтверждение-учётной-записи-мифи-и-номера-учебной-группы}{}

\hypertarget{tc-auth-004}{}

\textbf{Условия.} У А есть подтверждаемая учётная запись на университетском портале home.mephi.ru и номер группы в университетском профиле. Даны два независимых варианта. В первом у А ещё нет локальной записи, во втором в локальной записи уже есть личные решения. У Б отдельные университетский и локальный профили.

\textbf{Действие.} Для новой регистрации передать \texttt{provider\ =\ MEPHI} на маршрут регистрации. Для привязки сначала войти в существующий локальный профиль и обратиться к \texttt{POST /\allowbreak{}api/\allowbreak{}users/\allowbreak{}me/\allowbreak{}mephi-\allowbreak{}link}. В каждом варианте использовать полученные идентификатор обращения \texttt{auth\allowbreak{}Request\allowbreak{}Id} и адрес перехода \texttt{continue\allowbreak{}Url}. После подтверждения на портале завершить именно это обращение через \texttt{POST /\allowbreak{}api/\allowbreak{}auth/\allowbreak{}complete} в том же клиентском контексте. Затем выполнить вход по уже связанной университетской записи. Сопоставить номер и UUID группы с доверенным профилем и перечнем учебных групп.

\textbf{Ожидаемый результат.} Успешное завершение даёт 200 с токеном, идентификатором пользователя, UUID и номером подтверждённой группы. При привязке сохраняются исходный локальный UUID и история решений. Последующий вход через университетский портал возвращает того же пользователя. Для каждой группы в перечне различимы UUID, номер и участники, открытых паролей и хешей в ответе нет.

\textbf{Отказы.} В отдельных вариантах представить только номер группы строкой, совпадающий адрес электронной почты или один \texttt{auth\allowbreak{}Request\allowbreak{}Id} без подтверждения. Эти значения сами по себе не связывают профили и не назначают группу. Недействительное подтверждение даёт 401. При временной недоступности службы внешнего подтверждения ожидается 503 без создания записи или изменения прежней привязки. Повторная регистрация по уже связанной университетской записи или попытка связать её с другим локальным профилем даёт 409 без изменения прежней связи. Если у пользователя отсутствует подтверждённая группа, \texttt{group\allowbreak{}Number} и \texttt{groupId} равны \texttt{null}. Список шаблонов по запрошенной группе даёт 403, чужая запись отсутствует в личном расписании. Начало мероприятия без подтверждённой группы даёт 404 \texttt{ASSESSMENT\_\allowbreak{}NOT\_\allowbreak{}FOUND} без раскрытия назначения. Личные операции действительной локальной записи сохраняются. Ввести номер группы вручную для получения доступа нельзя.

\subsection{TC-AUTH-005. Назначение групп преподавателю и отзыв прав}\label{tc-auth-005-ux43dux430ux437ux43dux430ux447ux435ux43dux438ux435-ux433ux440ux443ux43fux43f-ux43fux440ux435ux43fux43eux434ux430ux432ux430ux442ux435ux43bux44e-ux438-ux43eux442ux437ux44bux432-ux43fux440ux430ux432}

\hypertarget{tc-auth-005-назначение-групп-преподавателю-и-отзыв-прав}{}

\hypertarget{tc-auth-005}{}

\textbf{Условия.} Администратор имеет роль \texttt{ADMIN}, преподаватель --- \texttt{TEACHER}, студент --- \texttt{USER}. Есть группы А и Б, личные сессии студентов обеих групп и черновик самого преподавателя. Преподавателю пока не назначены группы.

\textbf{Действие.} Под учётной записью администратора получить перечень преподавателей через \texttt{GET /\allowbreak{}api/\allowbreak{}admin/\allowbreak{}teachers}. Через \texttt{PUT /\allowbreak{}api/\allowbreak{}admin/\allowbreak{}teachers/\allowbreak{}\{teacher\allowbreak{}Id\}/\allowbreak{}groups} с массивом \texttt{groupIds} назначить преподавателю только группу А. Владелец черновика пробует опубликовать его для групп А и Б в отдельных вариантах. Для своей опубликованной работы группы А запросить результаты через \texttt{GET /\allowbreak{}api/\allowbreak{}assessments/\allowbreak{}\{assessment\allowbreak{}Id\}/\allowbreak{}results}. Администратор отзывает группу А. С тем же преподавательским токеном повторить запрос результатов и попытаться опубликовать новый черновик для этой группы. Отдельно запросить изменение назначений с правами только преподавателя и только студента. Администратор открывает общий перечень сессий.

\textbf{Ожидаемый результат.} Только администратор изменяет перечень групп преподавателя. Пользователи без роли \texttt{ADMIN} получают 403 при попытке назначения. Неизвестный UUID преподавателя или группы даёт 404 без частичного изменения прежнего списка. После назначения преподаватель публикует свои мероприятия для группы А и читает результаты собственных работ этой группы. При попытке получить результаты чужой работы или мероприятия соседней группы ожидается 404 без данных. Публикация для соседней группы отклоняется. Пользователь без роли \texttt{TEACHER} на маршруте результатов получает 403. После отзыва публикация для группы А отклоняется, чтение результатов ранее опубликованной работы даёт 404 без раскрытия данных. По маршруту общего обзора администратору доступны сессии студентов обеих групп.

\section{Шаблоны, контрольные мероприятия и модерация}\label{ux448ux430ux431ux43bux43eux43dux44b-ux43aux43eux43dux442ux440ux43eux43bux44cux43dux44bux435-ux43cux435ux440ux43eux43fux440ux438ux44fux442ux438ux44f-ux438-ux43cux43eux434ux435ux440ux430ux446ux438ux44f}

\hypertarget{шаблоны-контрольные-мероприятия-и-модерация}{}

\subsection{TC-CFG-001. Доступность и порядок шаблонов задач}\label{tc-cfg-001-ux434ux43eux441ux442ux443ux43fux43dux43eux441ux442ux44c-ux438-ux43fux43eux440ux44fux434ux43eux43a-ux448ux430ux431ux43bux43eux43dux43eux432-ux437ux430ux434ux430ux447}

\hypertarget{tc-cfg-001-доступность-и-порядок-шаблонов-задач}{}

\hypertarget{tc-cfg-001}{}

\textbf{Условия.} У студента роль \texttt{USER}, означающая право выполнять студенческие операции. Для группы А существуют включённый и выключенный шаблоны с разными именами. Во включённом шаблоне задана последовательность из трёх позиций, включая две задачи одного типа. Во второй группе есть одноимённый шаблон с иным составом.

\textbf{Действие.} С ролью \texttt{USER} запросить перечень для группы А, затем для другой группы. Повторить запрос списка от имени пользователя без группы. В перечне своей группы сопоставить возвращённые имена, последовательности типов и признаки включения \texttt{enabled}. Попытаться создать набор по выключенному шаблону. Затем создать его по включённому.

\textbf{Ожидаемый результат.} Студент видит оба шаблона своей группы с их именами и признаками включения. Для чужого \texttt{groupId} и для студента без группы запрос списка даёт 403 без выдачи шаблонов. При создании по недоступному или выключенному шаблону ожидается 404 \texttt{CONFIGURATION\_\allowbreak{}NOT\_\allowbreak{}FOUND}. По включённому выданы задачи в порядке шаблона группы А, две позиции одного типа имеют разные UUID. Данные одноимённого шаблона другой группы не подменяют состав созданной сессии.

\subsection{TC-CFG-002. Обзор модератора}\label{tc-cfg-002-ux43eux431ux437ux43eux440-ux43cux43eux434ux435ux440ux430ux442ux43eux440ux430}

\hypertarget{tc-cfg-002-обзор-модератора}{}

\hypertarget{tc-cfg-002}{}

\textbf{Условия.} Даны открытая сессия А и завершённая сессия Б из разных групп. Модератор имеет роль \texttt{ADMIN}, которая разрешает общий просмотр сессий. У обычного студента только роль \texttt{USER}.

\textbf{Действие.} Запросить обзор с обеих учётных записей. Сравнить принадлежность и состояния показанных сессий и задач с исходными данными, затем прочитать свои сессии от имени студентов.

\textbf{Ожидаемый результат.} Пользователь с ролью \texttt{ADMIN} получает обе сессии с состояниями задач. Для студента с ролью \texttt{USER} ожидается 403 без выдачи общего обзора. Для каждой доступной модератору записи различимы владелец и набор. Данные и результаты студентов после просмотра сохраняются.

\textbf{Роли, срок и ведомость преподавателя.} Дополнительно есть преподаватель с ролью \texttt{TEACHER}, собственная опубликованная работа в закреплённой группе и чужая работа. С этой ролью общий административный обзор даёт 403. Свои результаты через \texttt{/\allowbreak{}api/\allowbreak{}assessments/\allowbreak{}\{assessment\allowbreak{}Id\}/\allowbreak{}results} доступны; чужая работа, черновик или отзыв назначения группы дают одинаковый 404 без выдачи результатов. После отзыва назначения личные сессии и вердикты студентов сохраняются.

Для доступной ведомости сопоставить одно прохождение при чтении преподавателем, владельцем и администратором. Общий срок \texttt{D} --- 10:30:00, фоновый проход ещё не состоялся. При отсутствии ответа ровно в \texttt{D} и в 10:30:01 все три представления дают задаче \texttt{FAILED} и \texttt{result\ =\ null}. На независимом входе полный допустимый ответ принят в 10:29:59.900. Пока он обрабатывается, в \texttt{D} и в 10:30:01 задача имеет \texttt{NEW}; исходный предел \texttt{L} --- 10:30:59.900. Ровно в \texttt{L} без вердикта и другой своевременной обработки ожидается \texttt{FAILED} с \texttt{null}. Если вердикт сохранён до \texttt{L}, при позднейшем чтении сохраняются он и соответствующее состояние. Для незавершённого ответа с \texttt{L\ \textless{}\ D} таймаут оставляет задачу открытой до общего срока. Итог сессии вычисляется по состояниям всех её задач; эталоны и приватный контекст не выдаются. Это ожидаемые результаты описательных случаев, не запись выполненного прогона.

\subsection{TC-CFG-003. Состав контрольного мероприятия}\label{tc-cfg-003-ux441ux43eux441ux442ux430ux432-ux43aux43eux43dux442ux440ux43eux43bux44cux43dux43eux433ux43e-ux43cux435ux440ux43eux43fux440ux438ux44fux442ux438ux44f}

\hypertarget{tc-cfg-003-состав-контрольного-мероприятия}{}

\hypertarget{tc-cfg-003}{}

\textbf{Условия.} Преподаватель готовит контрольное мероприятие. Для него выбраны три позиции в заданном порядке --- ранцевая задача, задача Эль-Гамаля и ещё одна ранцевая задача. У двух студентов группы пока нет личных сессий по этому мероприятию.

\textbf{Действие.} Создать черновик через \texttt{POST /\allowbreak{}api/\allowbreak{}assessments}, указав поле \texttt{repeat\allowbreak{}Allowed}, которое задаёт правило повторного старта. Значение \texttt{true} разрешает повтор, \texttt{false} запрещает его. Записать его \texttt{assessment\allowbreak{}Id}, найти свой черновик в списке мероприятий и сопоставить каждую позицию с выбором преподавателя. Изменить порядок позиций и значение \texttt{repeat\allowbreak{}Allowed} в своём черновике, прочитать его повторно. При создании и при полной замене черновика методом \texttt{PUT} отдельно разобрать пропуск поля, значение \texttt{null} и строку вместо логического значения. Затем обратиться к личным сессиям обоих студентов.

\textbf{Ожидаемый результат.} По идентификатору найдено подготовленное мероприятие. Состояние \texttt{DRAFT} означает, что это черновик. Повторное чтение показывает тот же идентификатор, отредактированный порядок и выбранное значение \texttt{repeat\allowbreak{}Allowed}. Пропущенное поле и значение \texttt{null} допустимы для черновика и означают, что выбор ещё не сделан. Строка вместо логического значения даёт 400 без изменения ранее сохранённого черновика. Последовательность типов совпадает с заданной, в ней три позиции, включая две отдельные позиции одного типа. Само составление мероприятия не создаёт для студентов личных задач, решений и попыток.

\subsection{TC-CFG-004. Назначение мероприятия учебной группе}\label{tc-cfg-004-ux43dux430ux437ux43dux430ux447ux435ux43dux438ux435-ux43cux435ux440ux43eux43fux440ux438ux44fux442ux438ux44f-ux443ux447ux435ux431ux43dux43eux439-ux433ux440ux443ux43fux43fux435}

\hypertarget{tc-cfg-004-назначение-мероприятия-учебной-группе}{}

\hypertarget{tc-cfg-004}{}

\textbf{Условия.} Есть подготовленный черновик, две группы с разными UUID и выбранное преподавателем время открытия и закрытия. Известны идентификаторы мероприятия, выбранной группы и второй группы. В расписании уже есть другое мероприятие с тем же названием.

\textbf{Действие.} Опубликовать черновик для первой группы через \texttt{POST /\allowbreak{}api/\allowbreak{}assessments/\allowbreak{}\{assessment\allowbreak{}Id\}/\allowbreak{}publish}, передав UUID назначенной преподавателю группы в поле \texttt{groupId}, обязательное время открытия в \texttt{opensAt} и закрытия в \texttt{closesAt}. В независимых вариантах пропустить одну из границ, указать совпадающие границы или время закрытия раньше открытия. Сверить опубликованные данные с исходным составом, временем и выбранной группой. На отдельном черновике разобрать сбой записи при публикации.

\textbf{Ожидаемый результат.} По успешной публикации состояние становится \texttt{PUBLISHED} («опубликовано»). UUID мероприятия, его состав с редакциями вариантов, целевая группа, признак разрешения повтора и обе временные границы зафиксированы вместе. При отсутствии границы или нарушении строгого условия \texttt{opens\allowbreak{}At < closes\allowbreak{}At} публикация отклоняется, черновик сохраняется без строки в студенческом расписании. Для непустой работы с допустимой группой и окном отдельно рассмотреть два состояния черновика без выбора --- поле пропущено или равно \texttt{null}. При публикации ожидается 409 \texttt{ASSESSMENT\_\allowbreak{}RETAKE\_\allowbreak{}POLICY\_\allowbreak{}REQUIRED}. Работа остаётся \texttt{DRAFT} и скрыта в расписании. Если к тому же состав пуст, при корректных группе и окне первым возвращается 409 \texttt{ASSESSMENT\_\allowbreak{}EMPTY}. Непустую работу с явно выбранным \texttt{false} или \texttt{true} можно опубликовать, в \texttt{PUBLISHED} значение всегда логическое. Одноимённое мероприятие с другим UUID сохраняет собственную запись. Для второй группы назначение отсутствует. Если запись публикации прервалась, частичного назначения не возникает. Такой черновик недоступен студенту в расписании.

\subsection{TC-CFG-005. Автоматическое открытие, закрытие и повторный старт}\label{tc-cfg-005-ux430ux432ux442ux43eux43cux430ux442ux438ux447ux435ux441ux43aux43eux435-ux43eux442ux43aux440ux44bux442ux438ux435-ux437ux430ux43aux440ux44bux442ux438ux435-ux438-ux43fux43eux432ux442ux43eux440ux43dux44bux439-ux441ux442ux430ux440ux442}

\hypertarget{tc-cfg-005-автоматическое-открытие-закрытие-и-повторный-старт}{}

\hypertarget{tc-cfg-005}{}

\textbf{Условия.} Преподаватель с правом на группу опубликовал работу. У студентов А и Б подтверждены профили этой группы. Окно открыто с 10:00:00.000Z включительно до 11:00:00.000Z исключительно. Для независимых вариантов преподаватель подготовил два мероприятия. В одном черновике повторы разрешены (\texttt{repeat\allowbreak{}Allowed =\allowbreak{} true}), во втором запрещены (\texttt{false}). У студентов нет активных сессий.

\textbf{Границы окна.} При отсутствии прежних попыток запросить старт по UUID в параметре \texttt{assessment\allowbreak{}Id} за миллисекунду до открытия, точно в 10:00, за миллисекунду до закрытия и точно в 11:00. Повторить последний вариант ещё через миллисекунду. До открытия ожидается 409 \texttt{ASSESSMENT\_\allowbreak{}NOT\_\allowbreak{}OPEN}. В 10:00 окно открывается автоматически. В момент открытия и до 11:00 допустимый старт создаёт личную сессию. Ровно в 11:00 и позже возвращается 409 \texttt{ASSESSMENT\_\allowbreak{}CLOSED} без новой сессии. Студенту без предыдущего старта поле расписания \texttt{canStart} показывает разрешение на момент чтения. Если он открыл страницу до 11:00, но отправил запрос в 11:00, сервер отказывает по времени запроса.

\textbf{Повтор.} Внутри окна начать работу и ещё раз потребовать старт, пока первая сессия открыта (\texttt{NEW}). При обоих значениях \texttt{repeat\allowbreak{}Allowed} ожидается 409 \texttt{ACTIVE\_\allowbreak{}SESSION\_\allowbreak{}EXISTS} с UUID собственной действующей сессии. Завершить все задачи и повторить старт до 11:00. При разрешённых повторах выдана новая личная сессия с новым UUID. Прежняя сессия и её результат сохранены и относятся к тому же мероприятию. При запрещённых повторах приходит 409 \texttt{ASSESSMENT\_\allowbreak{}RETAKE\_\allowbreak{}DISABLED} с UUID прежней личной сессии в \texttt{details.\allowbreak{}session\allowbreak{}Id}, второй сессии нет. Ровно с 11:00 ожидается \texttt{ASSESSMENT\_\allowbreak{}CLOSED} и при разрешённых, и при запрещённых повторах, в том числе если ранее начатая попытка ещё ожидает своевременный вердикт. Время для видимой работы проверяется раньше состояния сессии и настройки повтора.

\textbf{Принадлежность и практика.} В открытом окне начинают А и Б. У каждого свои UUID сессии и задач, в поле \texttt{assessment\allowbreak{}Id} --- одно и то же мероприятие. В личном расписании \texttt{sessionId} указывает на активную сессию студента по этому мероприятию, а после завершения --- на его последнюю сессию по тому же мероприятию. UUID чужого старта отсутствует. На независимом входе создать сессию практики по доступному шаблону с параметром \texttt{configuration} и завершить задачу. Поле \texttt{assessment\allowbreak{}Id} равно \texttt{null}. Сессия и вердикт практики не образуют участия или оценки мероприятия. При передаче обоих параметров, а также при отсутствии обоих, ожидается 400 \texttt{INVALID\_\allowbreak{}REQUEST} без создания сессии. Сроки отправки решения по уже начатой работе разобраны в {\hyperlink{tc-ses-017}{TC-SES-017}}.

\subsection{TC-CFG-006. Авторское задание с текстовым ответом}\label{tc-cfg-006-ux430ux432ux442ux43eux440ux441ux43aux43eux435-ux437ux430ux434ux430ux43dux438ux435-ux441-ux442ux435ux43aux441ux442ux43eux432ux44bux43c-ux43eux442ux432ux435ux442ux43eux43c}

\hypertarget{tc-cfg-006-авторское-задание-с-текстовым-ответом}{}

\hypertarget{tc-cfg-006}{}

\textbf{Условия.} Преподаватель работает в форме создания задания \texttt{CUSTOM\_\allowbreak{}TEXT} --- вида задания с текстовым ответом. Он вводит две пары. В первом варианте условие --- «Назовите код первой группы», эталон --- « Код А ». Во втором условие --- «Назовите код второй группы», эталон --- «Код Б». Пробелы на краях первого эталона входят в исходные данные.

\textbf{Действие.} Сохранить варианты средствами формы, без изменения программы. В запросе создания авторского задания указывают название, а для каждого варианта передают условие в поле \texttt{statement} и закрытый правильный ответ в поле \texttt{expected\allowbreak{}Answer}. На независимых входах подать пустой массив вариантов, одну заполненную пару, затем массив из обеих заполненных пар. Для выдачи первого варианта заполнить текстовое поле «Ответ» значением «Код А » и передать его на проверку. На независимых открытых задачах с тем же вариантом передать «код А», «Код А» с двумя внутренними пробелами и «Код Б». Повторить последний ответ на задаче со вторым вариантом.

\textbf{Ожидаемый результат.} Пустой массив даёт 400 без создания задания. Одна заполненная пара допустима и всегда выбирается при выдаче. При двух или более парах выбор случайный и закрепляется по правилам {\hyperlink{tc-cfg-007}{TC-CFG-007}}. У каждой пары своё условие и свой закрытый эталон. При сравнении удаляются крайние пробелы у обеих строк. Для первого варианта ответ «Код А » верен. Изменение регистра, числа внутренних пробелов или подстановка эталона второй пары дают неверный предметный ответ. Для второго варианта «Код Б» верен. Каждый вердикт относится к той задаче и той паре, которые были выданы студенту.

\textbf{Неполная пара.} На независимых входах создания и полной замены собственного упражнения рассмотреть отсутствие statement, отсутствие expectedAnswer и эталон числового типа вместо строки. В каждом случае ожидается 400 без сохранения неполного варианта; при замене прежние полные пары сохраняются. Такая пара не появляется в банке и не участвует в публикации или выдаче. Затем сохранить одну полную пару и сопоставить её условие с соответствующим закрытым эталоном.

\subsection{TC-CFG-007. Закрепление выданного варианта}\label{tc-cfg-007-ux437ux430ux43aux440ux435ux43fux43bux435ux43dux438ux435-ux432ux44bux434ux430ux43dux43dux43eux433ux43e-ux432ux430ux440ux438ux430ux43dux442ux430}

\hypertarget{tc-cfg-007-закрепление-выданного-варианта}{}

\hypertarget{tc-cfg-007}{}

\textbf{Условия.} В опубликованное мероприятие входят две позиции авторского вида. У каждой позиции есть по меньшей мере две подготовленные пары с различимыми условиями и эталонами. В группе два студента. Открытие уже наступило.

\textbf{Действие.} После публикации преподаватель-составитель меняет название и пары вариантов через \texttt{PUT /\allowbreak{}api/\allowbreak{}custom-\allowbreak{}tasks/\allowbreak{}\{custom\allowbreak{}Task\allowbreak{}Id\}}. По списку собственных заданий сверить новое название (\texttt{title}) и список вариантов (\texttt{variants}). Эти изменения доступны для следующей публикации. При старте мероприятия каждым студентом проследить выбор варианта на каждую позицию его личной сессии. Для одной позиции на независимых исходных данных разобрать каждый возможный исход выбора. Записать UUID задач и выданные условия. Повторить получение сессии, перезапустить приложение и прочитать её снова. Сопоставить ответ, выданный студенту, со всеми полями варианта, доступным преподавателю. После решения прочитать личную историю.

\textbf{Ожидаемый результат.} Для каждой позиции ранее опубликованной работы используются варианты редакции, зафиксированной при её публикации. Если преподаватель включает то же авторское задание в новую публикацию после редактирования, в ней доступны обновлённые пары. Случайно выбранный вариант относится к конкретной позиции и задаче этого студента. Клиент получает только условие своей задачи. В выдаче при создании, при чтении, после перезапуска и в истории решённых задач отсутствуют закрытый эталон и чужие варианты. У каждой задачи сохраняются её UUID и условие, проверка использует соответствующий эталон. Назначение одинакового варианта двум студентам допустимо. Повторное чтение уже созданной личной сессии сохраняет прежний выбор.

\subsection{TC-CFG-008. Публикация и видимость в расписании}\label{tc-cfg-008-ux43fux443ux431ux43bux438ux43aux430ux446ux438ux44f-ux438-ux432ux438ux434ux438ux43cux43eux441ux442ux44c-ux432-ux440ux430ux441ux43fux438ux441ux430ux43dux438ux438}

\hypertarget{tc-cfg-008-публикация-и-видимость-в-расписании}{}

\hypertarget{tc-cfg-008}{}

\textbf{Условия.} Преподаватель подготовил черновик с заданиями разных типов и выбрал группу А для публикации. Дата и время открытия ещё не наступили. Есть участники групп А и Б, у каждого отсутствует личная сессия по назначению.

\textbf{Действие.} До публикации получить личное расписание под учётными записями студентов обеих групп через \texttt{GET /\allowbreak{}api/\allowbreak{}users/\allowbreak{}me/\allowbreak{}schedule}. Опубликовать мероприятие явным действием преподавателя с указанием группы и времени. Сразу после публикации, всё ещё до открытия, обновить расписание каждого студента и попытаться начать мероприятие. Сопоставить сведения строки расписания с назначением и прочитать их повторно.

\textbf{Ожидаемый результат.} Черновик отсутствует в студенческом расписании. Сразу после публикации в расписании А появляется запись о том же мероприятии с заданными временем открытия и временем закрытия. До личного старта запись не содержит UUID сессии другого участника. У студента Б запись отсутствует. Для чужого или неопубликованного мероприятия он получает 404 \texttt{ASSESSMENT\_\allowbreak{}NOT\_\allowbreak{}FOUND} как до открытия, так и после закрытия. Время и наличие чужих сессий остаются скрытыми. Для студента А до открытия в поле доступности старта указано \texttt{can\allowbreak{}Start =\allowbreak{} false}, а запрос старта даёт 409 \texttt{ASSESSMENT\_\allowbreak{}NOT\_\allowbreak{}OPEN}, задачи, тексты вариантов и эталоны не выдаются. Чтение строки расписания и раннее обращение за стартом сохраняют состояние назначения и число личных сессий.

\textbf{Граница публикации --- ноль и одно задание.} Создать для преподавателя отдельный черновик, передав название и \texttt{"items":\ {[}{]}}. Ожидается 201 и собственный \texttt{assessment\allowbreak{}Id}. Для него передать допустимые \texttt{groupId}, \texttt{opensAt} и \texttt{closesAt} в запросе публикации; обе границы заданы и \texttt{opens\allowbreak{}At < closes\allowbreak{}At}. Выбор \texttt{repeat\allowbreak{}Allowed} пока отложен. Ожидается 409 с кодом \texttt{ASSESSMENT\_\allowbreak{}EMPTY} и сообщением о необходимости добавить задание. Проверить по списку преподавателя, что тот же черновик находится в состоянии \texttt{DRAFT}. У студентов обеих групп он отсутствует в расписании.

Сохранить в этом черновике одну допустимую позицию, затем удалить её через \texttt{PUT /\allowbreak{}api/\allowbreak{}assessments/\allowbreak{}\{assessment\allowbreak{}Id\}}, передав название и пустой массив \texttt{items}. Ожидается 200 при сохранении и тот же отказ при публикации. Вернуть одну позицию, сохранив название и \texttt{repeat\allowbreak{}Allowed =\allowbreak{} null}, и повторить публикацию с прежними идентификатором, группой и обеими границами окна. Ожидается 409 \texttt{ASSESSMENT\_\allowbreak{}RETAKE\_\allowbreak{}POLICY\_\allowbreak{}REQUIRED}, тот же \texttt{DRAFT} и отсутствие записи расписания. Затем сохранить название, одну позицию и явное \texttt{repeat\allowbreak{}Allowed =\allowbreak{} false}. С теми же допустимыми параметрами назначения повторить публикацию. Ожидается 200 и состояние \texttt{PUBLISHED}. До открытия группа А уже видит одну запись с этим \texttt{assessment\allowbreak{}Id}, группа Б её не видит. Сама публикация не создаёт личных сессий.

Подготовить два черновика с одной позицией и явным \texttt{repeat\allowbreak{}Allowed =\allowbreak{} false} в каждом. В обоих запросах публикации передать допустимые группу и обе границы окна. Для каждого одновременно рассмотреть запрос публикации и запрос на удаление последнего задания. В описательной проверке задать порядок фиксации: для первого черновика сначала сохранить удаление, для второго --- публикацию. Если первым сохранено удаление, публикация должна завершиться отказом \texttt{ASSESSMENT\_\allowbreak{}EMPTY}. Если первой сохранена публикация, группа видит работу с одной позицией. Запрос на удаление получает 409 с кодом \texttt{ASSESSMENT\_\allowbreak{}ALREADY\_\allowbreak{}PUBLISHED}. Повторное чтение показывает исходную позицию.

\section{Ранцевые задачи}\label{ux440ux430ux43dux446ux435ux432ux44bux435-ux437ux430ux434ux430ux447ux438}

\hypertarget{ранцевые-задачи}{}

В учебной ранцевой задаче сообщение представляют последовательностью логических значений. Лёгкий и тяжёлый ранцы --- два связанных списка чисел. Генератор строит их по правилам выбранного варианта. Эталон связывает исходные величины, результат шифрования (шифртекст) и сообщение после расшифрования. Размерность задаёт число элементов ранца. Верификатор использует по одному учебному эталону для каждого из двух вариантов и {\href{../07\_Программные\_требования/DOC-06.3\_Ранцевые\_криптосистемы.pdf\#doc-063-ранцевые-криптосистемы}{спецификацию подсистемы}}. Ответы на разные варианты одной карточки поступают в независимые сессии до общего срока. Фиксация --- сохранение вердикта и связанных с ним состояний. Время фиксации предметного ответа в этих карточках укладывается в предел обработки.

\subsection{TC-KNAP-001. Выдача условия и сохранение выбранного варианта}\label{tc-knap-001-ux432ux44bux434ux430ux447ux430-ux443ux441ux43bux43eux432ux438ux44f-ux438-ux441ux43eux445ux440ux430ux43dux435ux43dux438ux435-ux432ux44bux431ux440ux430ux43dux43dux43eux433ux43e-ux432ux430ux440ux438ux430ux43dux442ux430}

\hypertarget{tc-knap-001-выдача-условия-и-сохранение-выбранного-варианта}{}

\hypertarget{tc-knap-001}{}

\textbf{Условия.} В шаблоне группы включена ранцевая задача. Внутреннее обозначение \texttt{Backpack} служит для выбора обработчика, который выдаёт её условие и принимает ответ. Для генератора подготовлены оба возможных исхода выбора варианта и примеры сообщений длиной 3, 4 и 5 элементов.

\textbf{Действие.} Проследить создание сессии отдельно при каждом исходе выбора. Сопоставить выданное условие с полным эталоном, посчитать длину сообщения, элементы со значением \texttt{false} и число раскрытых элементов лёгкого ранца. Для каждой позиции сверить сохраняемые данные с ответом на создание. После перезапуска прочитать выданное условие. Затем передать верное решение до срока и прочитать историю владельца.

\textbf{Ожидаемый результат.} Вызван обработчик нужного квалификатора. Случайный выбор даёт один из двух поддерживаемых вариантов. Доля ложных значений не меньше трети, поэтому при длине 3 их хотя бы одно, а при длине 4 или 5 --- хотя бы два. Открыты не более двух элементов лёгкого ранца. Вариант без выданного множителя также позволяет решить учебное задание по сохранённому условию. Параметры эталона согласуются с выданными фрагментами. Состояние, UUID, принадлежность, тип и условие сохранены до выдачи. После перезапуска возвращается то же условие. По модели хранения списковые величины записаны в текстовом представлении. Условие решённого задания доступно владельцу в истории.

\subsection{TC-KNAP-002. Форма ответа и запрет замены варианта}\label{tc-knap-002-ux444ux43eux440ux43cux430-ux43eux442ux432ux435ux442ux430-ux438-ux437ux430ux43fux440ux435ux442-ux437ux430ux43cux435ux43dux44b-ux432ux430ux440ux438ux430ux43dux442ux430}

\hypertarget{tc-knap-002-форма-ответа-и-запрет-замены-варианта}{}

\hypertarget{tc-knap-002}{}

\textbf{Условия.} Для каждого из двух вариантов дана открытая задача с допустимым решением. В решении присутствуют размерность, вариант, исходное и восстановленное сообщения, лёгкий и тяжёлый ранцы, множитель, модуль, обратный множитель и шифртекст. Модуль здесь --- число, по которому берут остаток. Множитель и обратный множитель --- величины учебного преобразования, произведение которых даёт остаток 1 по этому модулю. Поля и их типы перечислены в {\href{../07\_Программные\_требования/DOC-06.2\_Сессии\_задач\_и\_планировщик.pdf\#ранцевая-задача}{схеме ранцевого решения}}.

\textbf{Действие.} На независимых входах удалить обязательное поле, передать \texttt{module} JSON-числом вместо десятичной строки или указать вариант, отличный от сохранённого. Затем до срока отправить полный ответ с надлежащей формой.

\textbf{Ожидаемый результат.} В каждом из трёх отклонений клиент получает 400 \texttt{INVALID\_\allowbreak{}REQUEST}, предметного результата нет, попытка сохраняется. Исправленный запрос направляется проверяющему варианта выданной задачи.

\subsection{TC-KNAP-003. Правильность и первая ошибка}\label{tc-knap-003-ux43fux440ux430ux432ux438ux43bux44cux43dux43eux441ux442ux44c-ux438-ux43fux435ux440ux432ux430ux44f-ux43eux448ux438ux431ux43aux430}

\hypertarget{tc-knap-003-правильность-и-первая-ошибка}{}

\hypertarget{tc-knap-003}{}

\textbf{Условия.} Есть эталоны двух вариантов. Для каждого эталона подготовлены верный ответ и ответы, отличающиеся в одной проверяемой величине. Для диагностики дополнительно даны два несоответствия в одном ответе. Все ответы имеют правильную форму.

\textbf{Действие.} Сопоставить по этапам исходные данные, размерность, длины и элементы сообщений и ранцев, свойства лёгкого ранца выбранного варианта, модуль, взаимную простоту множителя с модулем (у них нет общего делителя больше единицы), тяжёлый ранец, обратный множитель, шифртекст и результат расшифрования. Для варианта на степенях проверить построение из степеней основания, для супервозрастающего --- положительность. Каждый очередной элемент лёгкого ранца супервозрастающего варианта должен быть больше суммы предшествующих. Проверить положительность элементов обоих ранцев. Проследить, как меняется диагностика при введении ошибки на каждом этапе, включая последний бит расшифрованного сообщения.

\textbf{Ожидаемый результат.} Верное решение каждого варианта получает 200 с \texttt{is\allowbreak{}Ok =\allowbreak{} true} и сохраняется как \texttt{SOLVED}. Ошибка в математической величине получает 200 с \texttt{is\allowbreak{}Ok =\allowbreak{} false}, задача переходит в \texttt{FAILED}. В случае с двумя ошибками сообщение указывает на первую неверную величину по порядку проверки. При проверке расшифрования по супервозрастающему ранцу применяется предусмотренный спецификацией жадный алгоритм. В этом варианте при разборе суммы элементы ранца рассматривают от большего к меньшему. Повторная подача до срока для любой итоговой задачи отклоняется с 409 \texttt{TASK\_\allowbreak{}FINALIZED}. Сохранённый результат остаётся доступным.

\subsection{TC-KNAP-004. Редакционная генерация}\label{tc-knap-004-ux440ux435ux434ux430ux43aux446ux438ux43eux43dux43dux430ux44f-ux433ux435ux43dux435ux440ux430ux446ux438ux44f}

\hypertarget{tc-knap-004-редакционная-генерация}{}

\hypertarget{tc-knap-004}{}

\textbf{Условия.} Редакционная генерация выдаёт образец задачи с полным решением для разбора. Для каждого варианта подготовлены два независимых запроса. В одном выбран только вариант, в другом также заданы согласованные размерность, сообщение, лёгкий ранец и множитель. В третьем, отрицательном, множитель имеет с модулем общий делитель больше единицы. Состав сохранённых задач и сессий известен до обращения.

\textbf{Действие.} Получить эталон по варианту, затем по явно указанным данным. Сравнить выданные величины с выбранным вариантом и переданными параметрами. Для запроса без данных проверить принадлежность автоматически выбранного основания заданному спецификацией набору простых чисел, если вариант требует выбора основания. Повторить запрос с несовместимым множителем.

\textbf{Ожидаемый результат.} Для каждого допустимого входа выбран соответствующий решатель и выдан полный эталон от исходного до расшифрованного сообщения. Переданные параметры сохранены в эталоне. Несовместимый множитель даёт отказ с указанием причины до построения результата. Ни один редакционный запрос не добавляет в хранилище условие, попытку или сессию.

\section{Эль-Гамаль и демонстрационные упражнения}\label{ux44dux43bux44c-ux433ux430ux43cux430ux43bux44c-ux438-ux434ux435ux43cux43eux43dux441ux442ux440ux430ux446ux438ux43eux43dux43dux44bux435-ux443ux43fux440ux430ux436ux43dux435ux43dux438ux44f}

\hypertarget{эль-гамаль-и-демонстрационные-упражнения}{}

\subsection{TC-ELG-001. Генерация условия и проверка решения}\label{tc-elg-001-ux433ux435ux43dux435ux440ux430ux446ux438ux44f-ux443ux441ux43bux43eux432ux438ux44f-ux438-ux43fux440ux43eux432ux435ux440ux43aux430-ux440ux435ux448ux435ux43dux438ux44f}

\hypertarget{tc-elg-001-генерация-условия-и-проверка-решения}{}

\hypertarget{tc-elg-001}{}

\textbf{Условия.} В составе сессии выдана задача по схеме Эль-Гамаля --- учебному примеру шифрования с открытым ключом. В условии \texttt{p} задаёт простой модуль, \texttt{g} --- основание, \texttt{x} --- секретный показатель учебного примера, \texttt{m} --- исходное сообщение. К этому условию подготовлены правильные значения открытого ключа, допустимого одноразового показателя, обеих частей шифртекста и восстановленного сообщения. Поля решения пояснены в {\href{../07\_Программные\_требования/DOC-06.2\_Сессии\_задач\_и\_планировщик.pdf\#эль-гамаль}{предметной схеме}}. Для каждого варианта ответа задача открыта.

\textbf{Действие.} Сопоставить выданное условие с эталоном и передать корректное решение. На независимых входах внести по одному математическому несоответствию на каждом шаге, сохраняя допустимую форму, затем одновременно исказить два шага с известным порядком. Отдельно изменить тип обязательного поля.

\textbf{Ожидаемый результат.} Верное решение возвращает положительный вердикт, неверный расчёт --- отрицательный с указанием первой неверной величины. Ошибка типа возвращает отказ формата до предметного вердикта. Результат принадлежит исходной задаче и читается вместе с ней.

\subsection{TC-ELG-002. Готовый разбор и условие для самостоятельного решения}\label{tc-elg-002-ux433ux43eux442ux43eux432ux44bux439-ux440ux430ux437ux431ux43eux440-ux438-ux443ux441ux43bux43eux432ux438ux435-ux434ux43bux44f-ux441ux430ux43cux43eux441ux442ux43eux44fux442ux435ux43bux44cux43dux43eux433ux43e-ux440ux435ux448ux435ux43dux438ux44f}

\hypertarget{tc-elg-002-готовый-разбор-и-условие-для-самостоятельного-решения}{}

\hypertarget{tc-elg-002}{}

\textbf{Условия.} Разрядность (\texttt{bitLength}) задаёт длину параметра в двоичных разрядах. Выбрана допустимая учебная разрядность. Для зарегистрированного студента заранее записаны состояние сессии, состав истории и доступность попытки.

\textbf{Действие.} Запросить редакционный пошаговый разбор и, отдельным обращением, самостоятельное задание с той же допустимой длиной. Сравнить состав двух результатов и повторно прочитать сессию и историю. Для каждой операции дополнительно передать нулевой и отрицательный параметр \texttt{bitLength}.

\textbf{Ожидаемый результат.} Разбор содержит условие и значения последовательных шагов до восстановленного сообщения. Самостоятельная задача выдаёт условие, по которому студент может выполнить эти шаги сам. Обе допустимые операции учитывают выбранную длину. Для \texttt{bit\allowbreak{}Length =\allowbreak{} 0} и \texttt{bit\allowbreak{}Length =\allowbreak{} -\allowbreak{}1} ожидается HTTP 400 без материала. Ни одна из операций не создаёт сохранённой сессионной задачи или вердикта и не расходует попытку.

\subsection{TC-BDR-001. Демонстрационное упражнение по готовому примеру}\label{tc-bdr-001-ux434ux435ux43cux43eux43dux441ux442ux440ux430ux446ux438ux43eux43dux43dux43eux435-ux443ux43fux440ux430ux436ux43dux435ux43dux438ux435-ux43fux43e-ux433ux43eux442ux43eux432ux43eux43cux443-ux43fux440ux438ux43cux435ux440ux443}

\hypertarget{tc-bdr-001-демонстрационное-упражнение-по-готовому-примеру}{}

\hypertarget{tc-bdr-001}{}

\textbf{Условия.} Преподаватель выбрал готовое учебное упражнение из заранее подготовленных примеров. Ему известны условие, опубликованный демонстрационный эталон и пояснение о том, какое допущение о конфиденциальности нарушено в этом примере. Конфиденциальность здесь означает, что содержание защищённого сообщения должно оставаться доступным только тем, кому оно предназначено. Состояние сессии, число доступных попыток и история студента известны.

\textbf{Действие.} Получить материалы упражнения, сопоставить представленное условие и итог с эталоном выбранного примера. По пояснению определить, какое допущение о конфиденциальности не выполнено и чем это отличается от неверно посчитанного значения в решении студента. Повторить просмотр. Отдельно рассмотреть недоступность готового материала. Сравнить сессию и историю до и после просмотра.

\textbf{Ожидаемый результат.} Студент получает материалы одного выбранного упражнения, результат соответствует его опубликованному эталону. В пояснении названо нарушенное допущение. Учебный результат иллюстрирует последствия этого допущения и отделён от диагностики вычислительной ошибки студента. При недоступном материале пользователь получает отказ без оценки своего решения. В обоих исходах сохраняются исходная сессия и история. Новая задача, запись результата или расход попытки отсутствуют.

\section{Интерфейс и эксплуатационные характеристики}\label{ux438ux43dux442ux435ux440ux444ux435ux439ux441-ux438-ux44dux43aux441ux43fux43bux443ux430ux442ux430ux446ux438ux43eux43dux43dux44bux435-ux445ux430ux440ux430ux43aux442ux435ux440ux438ux441ux442ux438ux43aux438}

\hypertarget{интерфейс-и-эксплуатационные-характеристики}{}

\subsection{TC-SYS-001. Общий контракт внешнего интерфейса}\label{tc-sys-001-ux43eux431ux449ux438ux439-ux43aux43eux43dux442ux440ux430ux43aux442-ux432ux43dux435ux448ux43dux435ux433ux43e-ux438ux43dux442ux435ux440ux444ux435ux439ux441ux430}

\hypertarget{tc-sys-001-общий-контракт-внешнего-интерфейса}{}

\hypertarget{tc-sys-001}{}

\textbf{Условия.} В справочнике API выбраны по одной операции с успехом и по одному представителю каждого применимого класса отказа. Ресурс --- объект, к которому клиент обращается по адресу, например сессия. Для запроса существующего объекта взят его UUID.

\textbf{Действие.} Сопоставить путь и метод каждой операции со справочником. Для операции с телом прочитать JSON при допустимом и недопустимом вводе. Последовательно рассмотреть неправильную форму, отсутствие учётных данных, нехватку роли и недоступный текущему пользователю объект. Проверить общий состав ошибки и дополнительные поля там, где они заданы конкретным контрактом.

\textbf{Ожидаемый результат.} Порт 8083 обозначает сетевую точку, на которой служба принимает запросы. Пути операций начинаются с \texttt{/api}. Для обмена структурированными данными используется JSON, у ресурсов сохраняются их UUID. Формат запроса, доступ и поиск ресурса получают разные статусы. Неверной форме соответствует 400, отсутствию подтверждённой личности --- 401, недостатку прав --- 403, отсутствию ресурса --- 404. Тело отказа соответствует общей схеме, а его код позволяет клиенту выбрать действие. Для ошибок сессий дополнительные договорённости и пограничные случаи приведены в карточках {\hyperlink{tc-ses-003}{TC-SES-003}}, {\hyperlink{tc-ses-005}{TC-SES-005}} и {\hyperlink{tc-ses-008}{TC-SES-008}}.

\subsection{TC-SYS-002. Передача секретов и параметров среды}\label{tc-sys-002-ux43fux435ux440ux435ux434ux430ux447ux430-ux441ux435ux43aux440ux435ux442ux43eux432-ux438-ux43fux430ux440ux430ux43cux435ux442ux440ux43eux432-ux441ux440ux435ux434ux44b}

\hypertarget{tc-sys-002-передача-секретов-и-параметров-среды}{}

\hypertarget{tc-sys-002}{}

\textbf{Условия.} Переменные окружения --- именованные настройки, которые приложение получает от среды запуска. Владелец развёртывания подготовил перечень имён внешних настроек и обезличенные примеры конфигурационных файлов. Для ключа подписи, срока JWT и параметров подключения к базе данных (БД) заданы контрольные значения, отдельные для двух сред.

\textbf{Действие.} Проследить источник каждого значения от запуска службы до её потребителя. Сравнить настройку двух сред при одной и той же поставке приложения. Проверить, что материалы поставки и примеры не содержат действующего ключа подписи и пароля БД.

\textbf{Ожидаемый результат.} Ключ подписи \texttt{JWT\_SECRET}, срок действия \texttt{ACCESS\_\allowbreak{}TOKEN\_\allowbreak{}LIFETIME} в днях и параметры доступа к БД поступают из переменных окружения. Смена среды не требует включать фактические секреты в исходные материалы. Для идентификации значений в записи проверки используются имена настроек, без раскрытия их содержимого.

\subsection{TC-SYS-003. Перезапуск сервиса и готовность базы}\label{tc-sys-003-ux43fux435ux440ux435ux437ux430ux43fux443ux441ux43a-ux441ux435ux440ux432ux438ux441ux430-ux438-ux433ux43eux442ux43eux432ux43dux43eux441ux442ux44c-ux431ux430ux437ux44b}

\hypertarget{tc-sys-003-перезапуск-сервиса-и-готовность-базы}{}

\hypertarget{tc-sys-003}{}

\textbf{Условия.} Данные хранятся в системе управления базами данных PostgreSQL 17. Средство Docker Compose задаёт порядок запуска служб в изолированных средах --- контейнерах. Постоянный том хранит файлы базы отдельно от срока жизни её контейнера. В PostgreSQL сохранены учебные пользователи, группа, шаблон набора, одна завершённая сессия с результатами и одна открытая. Последняя остаётся открытой до истечения общего срока. Для БД предусмотрен постоянный том.

\textbf{Действие.} Зафиксировать связи и состояния по UUID. Рассмотреть запуск, при котором база принимает соединения позднее процесса запуска служб. Перезапустить приложение, затем службы с подключением того же тома. Сравнить прочитанные объекты. В отдельном варианте общий срок открытого набора истекает во время перерыва, своевременных ответов в обработке нет.

\textbf{Ожидаемый результат.} В Docker Compose запуск прикладной службы зависит от успешной проверки готовности PostgreSQL. Сохранены учётные записи, членство, признаки включения шаблонов, UUID и условия задач, принадлежность сессий и подтверждённые результаты. До срока открытую сессию можно продолжить. Просроченные задачи без своевременного ответа завершаются по общему правилу. Подключение прежнего тома сохраняет данные между перезапусками служб.

\textbf{Перезапуск после публикации и до первого старта.} У преподавателя есть свой черновик, банк авторских заданий и отдельное опубликованное мероприятие для группы А. Записать UUID черновика, авторского задания, опубликованной работы, обе границы времени, явный запрет повтора (\texttt{repeat\allowbreak{}Allowed =\allowbreak{} false}) и редакцию вариантов. До первой личной сессии изменить банк и перезапустить приложение. Прочитать собственные черновики и банк преподавателя, а также расписание студента группы А. Начать мероприятие внутри опубликованного окна, завершить личную сессию и до закрытия попытаться начать снова.

\textbf{Ожидаемый результат.} У преподавателя сохранились черновик с прежним UUID и изменённый банк. Студент видит в расписании то же опубликованное назначение с той же группой, теми же временными границами и сохранённым \texttt{repeat\allowbreak{}Allowed =\allowbreak{} false}. После завершения первой попытки сервер возвращает 409 \texttt{ASSESSMENT\_\allowbreak{}RETAKE\_\allowbreak{}DISABLED}, новой сессии нет. Выдача в первой личной сессии опирается на редакцию, зафиксированную при публикации. Новые пары банка доступны для будущих публикаций и не подменяют исходное условие и закрытый эталон действующей работы.

\textbf{Повтор и сохранность нескольких прохождений.} На отдельном мероприятии явно разрешить повтор. Внутри окна завершить первое прохождение и создать второе. После перезапуска с тем же постоянным томом должны сохраниться разные UUID сессий, их связь с одним мероприятием, собственные задачи, закреплённые варианты и результаты первого прохождения. Расписание указывает на текущую сессию этой работы, а прежняя читается по своему UUID. Две одинаковые позиции типа остаются двумя задачами в прежнем порядке.

\textbf{Перезапуск с принятым ответом.} Рассмотреть остановку после сохранения своевременного полного тела, но до вердикта. После восстановления исходные \texttt{t\_in}, \texttt{L}, владелец, задача и тело не изменяются. До \texttt{L} обработку можно продолжить по исходному допуску. В \texttt{L} права записи уже нет; при \texttt{L\ \textless{}\ D} задача остаётся открытой, а после \texttt{D} без другой своевременной обработки закрывается без вердикта. Отдельно рассмотреть потерю соединения после совместной записи вердикта и состояний: читается сохранённый результат, повторной предметной проверки нет. Детальный порядок операций определён в DOC-06.2 и проверяется кейсами SES.

\subsection{TC-SYS-004. Время отклика пользовательских операций}\label{tc-sys-004-ux432ux440ux435ux43cux44f-ux43eux442ux43aux43bux438ux43aux430-ux43fux43eux43bux44cux437ux43eux432ux430ux442ux435ux43bux44cux441ux43aux438ux445-ux43eux43fux435ux440ux430ux446ux438ux439}

\hypertarget{tc-sys-004-время-отклика-пользовательских-операций}{}

\hypertarget{tc-sys-004}{}

\textbf{Условия.} Стенд --- среда, на которой оценивают работу системы. Перед измерением он прогрет. Его службы уже обрабатывали характерные для проверки запросы. Одновременно работают 30 пользователей с фиксированным составом операций и данных. К обычным запросам относятся регистрация, вход, чтение групп, шаблонов и личных данных, а также отправка решения со штатным проверяющим. Создание сессии с формированием задач относится к классу генерации.

\textbf{Действие.} Взять не менее 100 штатных наблюдений этого класса при неизменной нагрузке. Для каждого записать интервал от отправки запроса клиентом до получения полного ответа. Обычное ожидание очереди входит в интервал. Операции с искусственными отказами или специально внесёнными задержками анализируют отдельно. Упорядочить полученные длительности по возрастанию. При числе наблюдений N взять порядковый номер, равный округлённому вверх произведению 0,95 на N, считая первый элемент номером 1.

\textbf{Ожидаемый результат.} Выбранное значение, 95-й процентиль (p95), не превышает 2 секунд. При 100 наблюдениях сравнивают 95-е значение. Ровно 2 секунды удовлетворяют критерию, 2,001 секунды дают отклонение. Условия прогрева, нагрузку, число наблюдений и границы каждого интервала записывают вместе с выводом. Предел обработки одного ответа рассматривается отдельно в {\hyperlink{tc-ses-011}{TC-SES-011}}.

\subsection{TC-SYS-005. Время генерации учебного материала}\label{tc-sys-005-ux432ux440ux435ux43cux44f-ux433ux435ux43dux435ux440ux430ux446ux438ux438-ux443ux447ux435ux431ux43dux43eux433ux43e-ux43cux430ux442ux435ux440ux438ux430ux43bux430}

\hypertarget{tc-sys-005-время-генерации-учебного-материала}{}

\hypertarget{tc-sys-005}{}

\textbf{Условия.} На том же прогретом стенде работают 30 одновременных пользователей. В класс генерации входят создание сессии с формированием задач и учебные запросы генерации. Согласованы состав входов, разрядность учебных параметров там, где она задаётся, и размер выдачи.

\textbf{Действие.} Собрать отдельную выборку не менее чем из 100 штатных запросов генерации. Для каждого измерить время от отправки клиентом запроса до получения полного условия, набора или редакционного разбора, включая обычное ожидание. Зафиксировать условия и состав выборки. Обычные запросы из {\hyperlink{tc-sys-004}{TC-SYS-004}}, искусственные отказы и специально внесённые задержки в эту выборку не включать. По той же формуле выбрать p95 в этой выборке.

\textbf{Ожидаемый результат.} Для класса генерации p95 не превышает 5 секунд. При 100 запросах ровно 5 секунд на 95-й позиции соответствуют требованию. Значение 5,001 секунды на той же позиции превышает предел. Быстрые запросы на чтение не меняют оценку генерации.

\subsection{TC-SYS-006. Подключаемый вид задания}\label{tc-sys-006-ux43fux43eux434ux43aux43bux44eux447ux430ux435ux43cux44bux439-ux432ux438ux434-ux437ux430ux434ux430ux43dux438ux44f}

\hypertarget{tc-sys-006-подключаемый-вид-задания}{}

\hypertarget{tc-sys-006}{}

\textbf{Условия.} Модуль задания --- часть программы, которая для своего вида выдаёт условие и проверяет решение. Каталог перечисляет зарегистрированные виды, чтобы преподаватель мог выбрать их при составлении работы. В нём присутствуют ранцевая задача и Эль-Гамаль. Для разбора подключения подготовлен модуль «Сложение». У контрольного примера условие «2 + 3», верный ответ 5 и неверный ответ 4. Описание модуля задаёт выдачу условия и проверку ответа через интерфейс предметного вида. Для текстового задания преподавателя подготовлены два известных проверяющему варианта с разными эталонами.

\textbf{Действие.} Проверить регистрацию типа с устойчивым обозначением, выдачей условия и закрытых данных проверки, схемой допустимых полей ответа и предметной проверкой. Найти его через \texttt{GET /\allowbreak{}api/\allowbreak{}task-\allowbreak{}types}, включить рядом с двумя имеющимися видами в смешанное мероприятие и проследить выдачу условий. Для каждого вида передать допустимые верный и неверный ответы в независимых личных сессиях. Сопоставить путь регистрации и обращения к модулю со слоями сессий и планировщика. Открыть каталог типов от имени студента и прочитать всю запись текстового типа, включая вложенные поля схемы ответа. Сопоставить содержимое с двумя подготовленными вариантами.

\textbf{Ожидаемый результат.} Каждая позиция сохраняет свой тип. По одному интерфейсу подсистема получает условие у выбранного вида и передаёт ему на проверку выданное условие и ответ. Расширение каталога не требует изменения алгоритмов создания сессии, приёма решения и обработки срока. Для учебного примера результаты 5 и 4 дают соответственно положительный и отрицательный вердикты. В каталоге текстовый тип описывает одно поле для ввода ответа. Ни перечень допустимых значений, ни другая часть публичной записи не раскрывают варианты преподавателя, их эталоны или закрытые данные проверки. Встроенные виды по-прежнему получают свои данные. Правила принадлежности, срока и фиксации результата применяются одинаково.

\section{Связи с требованиями}\label{ux441ux432ux44fux437ux438-ux441-ux442ux440ux435ux431ux43eux432ux430ux43dux438ux44fux43cux438}

\hypertarget{связи-с-требованиями}{}

Таблица предназначена для переноса связей в DOC-09. Для сессий и общей интеграции полное соответствие карточкам TC-SES приведено вместе с их определениями в {\hyperlink{doc-10-верификация-сессий-задач-и-планировщика}{главе о сессиях}}. Графа «Программное основание» указывает только те SWR, для которых в комплекте имеется текст требования.

{\def\LTcaptype{none} % do not increment counter
\begin{longtable}[]{@{}
  >{\raggedright\arraybackslash}p{(\linewidth - 4\tabcolsep) * \real{0.2521}}
  >{\raggedright\arraybackslash}p{(\linewidth - 4\tabcolsep) * \real{0.3530}}
  >{\raggedright\arraybackslash}p{(\linewidth - 4\tabcolsep) * \real{0.3950}}@{}}
\toprule\noalign{}
\begin{minipage}[b]{\linewidth}\raggedright
Карточка
\end{minipage} & \begin{minipage}[b]{\linewidth}\raggedright
Системное основание
\end{minipage} & \begin{minipage}[b]{\linewidth}\raggedright
Программное основание
\end{minipage} \\
\midrule\noalign{}
\endhead
\bottomrule\noalign{}
\endlastfoot
{\hyperlink{tc-auth-001}{TC-AUTH-001}} & {\href{../04\_Системные\_требования/DOC-05\_Системные\_требования.pdf\#sr-auth-001}{SR-AUTH-001}} & {\href{../07\_Программные\_требования/DOC-06\_Программные\_требования.pdf\#swr-auth-001}{SWR-AUTH-001}} \\
{\hyperlink{tc-auth-002}{TC-AUTH-002}} & {\href{../04\_Системные\_требования/DOC-05\_Системные\_требования.pdf\#sr-auth-002}{SR-AUTH-002}}, {\href{../04\_Системные\_требования/DOC-05\_Системные\_требования.pdf\#sr-auth-003}{SR-AUTH-003}}, {\href{../04\_Системные\_требования/DOC-05\_Системные\_требования.pdf\#sr-auth-008}{SR-AUTH-008}} & {\href{../07\_Программные\_требования/DOC-06\_Программные\_требования.pdf\#swr-auth-002}{SWR-AUTH-002}}, {\href{../07\_Программные\_требования/DOC-06\_Программные\_требования.pdf\#swr-auth-003}{SWR-AUTH-003}}, {\href{../07\_Программные\_требования/DOC-06\_Программные\_требования.pdf\#swr-auth-008}{SWR-AUTH-008}}, {\href{../07\_Программные\_требования/DOC-06\_Программные\_требования.pdf\#swr-sys-001}{SWR-SYS-001}} \\
{\hyperlink{tc-auth-003}{TC-AUTH-003}} & {\href{../04\_Системные\_требования/DOC-05\_Системные\_требования.pdf\#sr-auth-004}{SR-AUTH-004}}, {\href{../04\_Системные\_требования/DOC-05\_Системные\_требования.pdf\#sr-auth-007}{SR-AUTH-007}}, {\href{../04\_Системные\_требования/DOC-05\_Системные\_требования.pdf\#sr-sys-001}{SR-SYS-001}} & {\href{../07\_Программные\_требования/DOC-06\_Программные\_требования.pdf\#swr-auth-004}{SWR-AUTH-004}}, {\href{../07\_Программные\_требования/DOC-06\_Программные\_требования.pdf\#swr-auth-007}{SWR-AUTH-007}}, {\href{../07\_Программные\_требования/DOC-06\_Программные\_требования.pdf\#swr-sys-001}{SWR-SYS-001}}, {\href{../07\_Программные\_требования/DOC-06\_Программные\_требования.pdf\#swr-sys-002}{SWR-SYS-002}} \\
{\hyperlink{tc-auth-004}{TC-AUTH-004}} & {\href{../04\_Системные\_требования/DOC-05\_Системные\_требования.pdf\#sr-auth-005}{SR-AUTH-005}}, {\href{../04\_Системные\_требования/DOC-05\_Системные\_требования.pdf\#sr-auth-006}{SR-AUTH-006}}, {\href{../04\_Системные\_требования/DOC-05\_Системные\_требования.pdf\#sr-ses-001}{SR-SES-001}} & {\href{../07\_Программные\_требования/DOC-06\_Программные\_требования.pdf\#swr-auth-005}{SWR-AUTH-005}}, {\href{../07\_Программные\_требования/DOC-06\_Программные\_требования.pdf\#swr-auth-006}{SWR-AUTH-006}}, {\href{../07\_Программные\_требования/DOC-06.2\_Сессии\_задач\_и\_планировщик.pdf\#swr-ses-002}{SWR-SES-002}} \\
{\hyperlink{tc-auth-005}{TC-AUTH-005}} & {\href{../04\_Системные\_требования/DOC-05\_Системные\_требования.pdf\#sr-auth-003}{SR-AUTH-003}}, {\href{../04\_Системные\_требования/DOC-05\_Системные\_требования.pdf\#sr-auth-009}{SR-AUTH-009}}, {\href{../04\_Системные\_требования/DOC-05\_Системные\_требования.pdf\#sr-cfg-002}{SR-CFG-002}}, {\href{../04\_Системные\_требования/DOC-05\_Системные\_требования.pdf\#sr-cfg-004}{SR-CFG-004}} & {\href{../07\_Программные\_требования/DOC-06\_Программные\_требования.pdf\#swr-auth-003}{SWR-AUTH-003}}, {\href{../07\_Программные\_требования/DOC-06\_Программные\_требования.pdf\#swr-auth-009}{SWR-AUTH-009}}, {\href{../07\_Программные\_требования/DOC-06\_Программные\_требования.pdf\#swr-cfg-002}{SWR-CFG-002}}, {\href{../07\_Программные\_требования/DOC-06\_Программные\_требования.pdf\#swr-cfg-004}{SWR-CFG-004}} \\
{\hyperlink{tc-cfg-001}{TC-CFG-001}} & {\href{../04\_Системные\_требования/DOC-05\_Системные\_требования.pdf\#sr-cfg-001}{SR-CFG-001}}, {\href{../04\_Системные\_требования/DOC-05\_Системные\_требования.pdf\#sr-sys-003}{SR-SYS-003}}, {\href{../04\_Системные\_требования/DOC-05\_Системные\_требования.pdf\#sr-ses-001}{SR-SES-001}} & {\href{../07\_Программные\_требования/DOC-06\_Программные\_требования.pdf\#swr-cfg-001}{SWR-CFG-001}}, {\href{../07\_Программные\_требования/DOC-06.2\_Сессии\_задач\_и\_планировщик.pdf\#swr-ses-002}{SWR-SES-002}}, {\href{../07\_Программные\_требования/DOC-06.2\_Сессии\_задач\_и\_планировщик.pdf\#swr-ses-004}{SWR-SES-004}}, {\href{../07\_Программные\_требования/DOC-06\_Программные\_требования.pdf\#swr-sys-003}{SWR-SYS-003}} \\
{\hyperlink{tc-cfg-002}{TC-CFG-002}} & {\href{../04\_Системные\_требования/DOC-05\_Системные\_требования.pdf\#sr-cfg-002}{SR-CFG-002}} & {\href{../07\_Программные\_требования/DOC-06\_Программные\_требования.pdf\#swr-cfg-002}{SWR-CFG-002}} \\
{\hyperlink{tc-cfg-003}{TC-CFG-003}} & {\href{../04\_Системные\_требования/DOC-05\_Системные\_требования.pdf\#sr-cfg-003}{SR-CFG-003}} & {\href{../07\_Программные\_требования/DOC-06\_Программные\_требования.pdf\#swr-cfg-003}{SWR-CFG-003}} \\
{\hyperlink{tc-cfg-004}{TC-CFG-004}} & {\href{../04\_Системные\_требования/DOC-05\_Системные\_требования.pdf\#sr-cfg-004}{SR-CFG-004}} & {\href{../07\_Программные\_требования/DOC-06\_Программные\_требования.pdf\#swr-cfg-004}{SWR-CFG-004}} \\
{\hyperlink{tc-cfg-005}{TC-CFG-005}} & {\href{../04\_Системные\_требования/DOC-05\_Системные\_требования.pdf\#sr-cfg-005}{SR-CFG-005}}, {\href{../04\_Системные\_требования/DOC-05\_Системные\_требования.pdf\#sr-ses-001}{SR-SES-001}}, {\href{../04\_Системные\_требования/DOC-05\_Системные\_требования.pdf\#sr-sys-002}{SR-SYS-002}} & {\href{../07\_Программные\_требования/DOC-06\_Программные\_требования.pdf\#swr-cfg-005}{SWR-CFG-005}}, {\href{../07\_Программные\_требования/DOC-06.2\_Сессии\_задач\_и\_планировщик.pdf\#swr-ses-004}{SWR-SES-004}}, {\href{../07\_Программные\_требования/DOC-06\_Программные\_требования.pdf\#swr-sys-002}{SWR-SYS-002}} \\
{\hyperlink{tc-cfg-006}{TC-CFG-006}} & {\href{../04\_Системные\_требования/DOC-05\_Системные\_требования.pdf\#sr-cfg-006}{SR-CFG-006}} & {\href{../07\_Программные\_требования/DOC-06\_Программные\_требования.pdf\#swr-cfg-006}{SWR-CFG-006}} \\
{\hyperlink{tc-cfg-007}{TC-CFG-007}} & {\href{../04\_Системные\_требования/DOC-05\_Системные\_требования.pdf\#sr-cfg-007}{SR-CFG-007}}, {\href{../04\_Системные\_требования/DOC-05\_Системные\_требования.pdf\#sr-sys-007}{SR-SYS-007}} & {\href{../07\_Программные\_требования/DOC-06\_Программные\_требования.pdf\#swr-cfg-007}{SWR-CFG-007}}, {\href{../07\_Программные\_требования/DOC-06.2\_Сессии\_задач\_и\_планировщик.pdf\#swr-ses-019}{SWR-SES-019}} \\
{\hyperlink{tc-cfg-008}{TC-CFG-008}} & {\href{../04\_Системные\_требования/DOC-05\_Системные\_требования.pdf\#sr-cfg-008}{SR-CFG-008}} & {\href{../07\_Программные\_требования/DOC-06\_Программные\_требования.pdf\#swr-cfg-008}{SWR-CFG-008}} \\
{\hyperlink{tc-knap-001}{TC-KNAP-001}} & {\href{../04\_Системные\_требования/DOC-05\_Системные\_требования.pdf\#sr-knap-001}{SR-KNAP-001}}, {\href{../04\_Системные\_требования/DOC-05\_Системные\_требования.pdf\#sr-knap-002}{SR-KNAP-002}}, {\href{../04\_Системные\_требования/DOC-05\_Системные\_требования.pdf\#sr-knap-004}{SR-KNAP-004}}, {\href{../04\_Системные\_требования/DOC-05\_Системные\_требования.pdf\#sr-knap-006}{SR-KNAP-006}}, {\href{../04\_Системные\_требования/DOC-05\_Системные\_требования.pdf\#sr-sys-007}{SR-SYS-007}} & {\href{../07\_Программные\_требования/DOC-06.3\_Ранцевые\_криптосистемы.pdf\#swr-knap-001}{SWR-KNAP-001}}, {\href{../07\_Программные\_требования/DOC-06.3\_Ранцевые\_криптосистемы.pdf\#swr-knap-002}{SWR-KNAP-002}}, {\href{../07\_Программные\_требования/DOC-06.3\_Ранцевые\_криптосистемы.pdf\#swr-knap-003}{SWR-KNAP-003}}, {\href{../07\_Программные\_требования/DOC-06.3\_Ранцевые\_криптосистемы.pdf\#swr-knap-004}{SWR-KNAP-004}}, {\href{../07\_Программные\_требования/DOC-06.3\_Ранцевые\_криптосистемы.pdf\#swr-knap-005}{SWR-KNAP-005}}, {\href{../07\_Программные\_требования/DOC-06.3\_Ранцевые\_криптосистемы.pdf\#swr-knap-006}{SWR-KNAP-006}}, {\href{../07\_Программные\_требования/DOC-06.3\_Ранцевые\_криптосистемы.pdf\#swr-knap-007}{SWR-KNAP-007}}, {\href{../07\_Программные\_требования/DOC-06.3\_Ранцевые\_криптосистемы.pdf\#swr-knap-016}{SWR-KNAP-016}} \\
{\hyperlink{tc-knap-002}{TC-KNAP-002}} & {\href{../04\_Системные\_требования/DOC-05\_Системные\_требования.pdf\#sr-knap-003}{SR-KNAP-003}}, {\href{../04\_Системные\_требования/DOC-05\_Системные\_требования.pdf\#sr-ses-009}{SR-SES-009}} & {\href{../07\_Программные\_требования/DOC-06.3\_Ранцевые\_криптосистемы.pdf\#swr-knap-008}{SWR-KNAP-008}}, {\href{../07\_Программные\_требования/DOC-06.2\_Сессии\_задач\_и\_планировщик.pdf\#swr-ses-009}{SWR-SES-009}} \\
{\hyperlink{tc-knap-003}{TC-KNAP-003}} & {\href{../04\_Системные\_требования/DOC-05\_Системные\_требования.pdf\#sr-knap-004}{SR-KNAP-004}}, {\href{../04\_Системные\_требования/DOC-05\_Системные\_требования.pdf\#sr-knap-005}{SR-KNAP-005}}, {\href{../04\_Системные\_требования/DOC-05\_Системные\_требования.pdf\#sr-knap-006}{SR-KNAP-006}}, {\href{../04\_Системные\_требования/DOC-05\_Системные\_требования.pdf\#sr-knap-007}{SR-KNAP-007}} & {\href{../07\_Программные\_требования/DOC-06.3\_Ранцевые\_криптосистемы.pdf\#swr-knap-005}{SWR-KNAP-005}}, {\href{../07\_Программные\_требования/DOC-06.3\_Ранцевые\_криптосистемы.pdf\#swr-knap-006}{SWR-KNAP-006}}, {\href{../07\_Программные\_требования/DOC-06.3\_Ранцевые\_криптосистемы.pdf\#swr-knap-009}{SWR-KNAP-009}}, {\href{../07\_Программные\_требования/DOC-06.3\_Ранцевые\_криптосистемы.pdf\#swr-knap-010}{SWR-KNAP-010}}, {\href{../07\_Программные\_требования/DOC-06.3\_Ранцевые\_криптосистемы.pdf\#swr-knap-011}{SWR-KNAP-011}} \\
{\hyperlink{tc-knap-004}{TC-KNAP-004}} & {\href{../04\_Системные\_требования/DOC-05\_Системные\_требования.pdf\#sr-knap-008}{SR-KNAP-008}}, {\href{../04\_Системные\_требования/DOC-05\_Системные\_требования.pdf\#sr-knap-009}{SR-KNAP-009}}, {\href{../04\_Системные\_требования/DOC-05\_Системные\_требования.pdf\#sr-knap-010}{SR-KNAP-010}} & {\href{../07\_Программные\_требования/DOC-06.3\_Ранцевые\_криптосистемы.pdf\#swr-knap-012}{SWR-KNAP-012}}, {\href{../07\_Программные\_требования/DOC-06.3\_Ранцевые\_криптосистемы.pdf\#swr-knap-013}{SWR-KNAP-013}}, {\href{../07\_Программные\_требования/DOC-06.3\_Ранцевые\_криптосистемы.pdf\#swr-knap-014}{SWR-KNAP-014}}, {\href{../07\_Программные\_требования/DOC-06.3\_Ранцевые\_криптосистемы.pdf\#swr-knap-015}{SWR-KNAP-015}} \\
{\hyperlink{tc-elg-001}{TC-ELG-001}} & {\href{../04\_Системные\_требования/DOC-05\_Системные\_требования.pdf\#sr-elg-001}{SR-ELG-001}}, {\href{../04\_Системные\_требования/DOC-05\_Системные\_требования.pdf\#sr-ses-004}{SR-SES-004}}, {\href{../04\_Системные\_требования/DOC-05\_Системные\_требования.pdf\#sr-ses-005}{SR-SES-005}} & {\href{../07\_Программные\_требования/DOC-06\_Программные\_требования.pdf\#swr-elg-001}{SWR-ELG-001}}, {\href{../07\_Программные\_требования/DOC-06.2\_Сессии\_задач\_и\_планировщик.pdf\#swr-ses-009}{SWR-SES-009}}, {\href{../07\_Программные\_требования/DOC-06.2\_Сессии\_задач\_и\_планировщик.pdf\#swr-ses-010}{SWR-SES-010}}, {\href{../07\_Программные\_требования/DOC-06.2\_Сессии\_задач\_и\_планировщик.pdf\#swr-ses-011}{SWR-SES-011}} \\
{\hyperlink{tc-elg-002}{TC-ELG-002}} & {\href{../04\_Системные\_требования/DOC-05\_Системные\_требования.pdf\#sr-elg-002}{SR-ELG-002}} & {\href{../07\_Программные\_требования/DOC-06\_Программные\_требования.pdf\#swr-elg-002}{SWR-ELG-002}} \\
{\hyperlink{tc-bdr-001}{TC-BDR-001}} & {\href{../04\_Системные\_требования/DOC-05\_Системные\_требования.pdf\#sr-bdr-001}{SR-BDR-001}} & {\href{../07\_Программные\_требования/DOC-06\_Программные\_требования.pdf\#swr-bdr-001}{SWR-BDR-001}} \\
{\hyperlink{tc-sys-001}{TC-SYS-001}} & {\href{../04\_Системные\_требования/DOC-05\_Системные\_требования.pdf\#sr-sys-008}{SR-SYS-008}}, {\href{../04\_Системные\_требования/DOC-05\_Системные\_требования.pdf\#sr-sys-009}{SR-SYS-009}}, {\href{../04\_Системные\_требования/DOC-05\_Системные\_требования.pdf\#sr-sys-010}{SR-SYS-010}} & {\href{../07\_Программные\_требования/DOC-06\_Программные\_требования.pdf\#swr-sys-007}{SWR-SYS-007}}, {\href{../07\_Программные\_требования/DOC-06\_Программные\_требования.pdf\#swr-sys-009}{SWR-SYS-009}}, {\href{../07\_Программные\_требования/DOC-06.2\_Сессии\_задач\_и\_планировщик.pdf\#swr-ses-022}{SWR-SES-022}} \\
{\hyperlink{tc-sys-002}{TC-SYS-002}} & {\href{../04\_Системные\_требования/DOC-05\_Системные\_требования.pdf\#sr-sys-011}{SR-SYS-011}} & {\href{../07\_Программные\_требования/DOC-06\_Программные\_требования.pdf\#swr-sys-010}{SWR-SYS-010}} \\
{\hyperlink{tc-sys-003}{TC-SYS-003}} & {\href{../04\_Системные\_требования/DOC-05\_Системные\_требования.pdf\#sr-sys-007}{SR-SYS-007}}, {\href{../04\_Системные\_требования/DOC-05\_Системные\_требования.pdf\#sr-sys-012}{SR-SYS-012}}, {\href{../04\_Системные\_требования/DOC-05\_Системные\_требования.pdf\#sr-sys-013}{SR-SYS-013}}, {\href{../04\_Системные\_требования/DOC-05\_Системные\_требования.pdf\#sr-cfg-004}{SR-CFG-004}}, {\href{../04\_Системные\_требования/DOC-05\_Системные\_требования.pdf\#sr-cfg-007}{SR-CFG-007}}, {\href{../04\_Системные\_требования/DOC-05\_Системные\_требования.pdf\#sr-cfg-008}{SR-CFG-008}} & {\href{../07\_Программные\_требования/DOC-06\_Программные\_требования.pdf\#swr-sys-008}{SWR-SYS-008}}, {\href{../07\_Программные\_требования/DOC-06\_Программные\_требования.pdf\#swr-sys-011}{SWR-SYS-011}}, {\href{../07\_Программные\_требования/DOC-06\_Программные\_требования.pdf\#swr-sys-012}{SWR-SYS-012}}, {\href{../07\_Программные\_требования/DOC-06.2\_Сессии\_задач\_и\_планировщик.pdf\#swr-ses-019}{SWR-SES-019}}, {\href{../07\_Программные\_требования/DOC-06\_Программные\_требования.pdf\#swr-cfg-004}{SWR-CFG-004}}, {\href{../07\_Программные\_требования/DOC-06\_Программные\_требования.pdf\#swr-cfg-007}{SWR-CFG-007}}, {\href{../07\_Программные\_требования/DOC-06\_Программные\_требования.pdf\#swr-cfg-008}{SWR-CFG-008}} \\
{\hyperlink{tc-sys-004}{TC-SYS-004}} & {\href{../04\_Системные\_требования/DOC-05\_Системные\_требования.pdf\#sr-sys-014}{SR-SYS-014}} & {\href{../07\_Программные\_требования/DOC-06\_Программные\_требования.pdf\#swr-sys-013}{SWR-SYS-013}} \\
{\hyperlink{tc-sys-005}{TC-SYS-005}} & {\href{../04\_Системные\_требования/DOC-05\_Системные\_требования.pdf\#sr-sys-015}{SR-SYS-015}} & {\href{../07\_Программные\_требования/DOC-06\_Программные\_требования.pdf\#swr-sys-014}{SWR-SYS-014}} \\
{\hyperlink{tc-sys-006}{TC-SYS-006}} & {\href{../04\_Системные\_требования/DOC-05\_Системные\_требования.pdf\#sr-sys-016}{SR-SYS-016}} & {\href{../07\_Программные\_требования/DOC-06\_Программные\_требования.pdf\#swr-sys-015}{SWR-SYS-015}} \\
\end{longtable}
}

\section{История редакций}\label{ux438ux441ux442ux43eux440ux438ux44f-ux440ux435ux434ux430ux43aux446ux438ux439}

\hypertarget{история-редакций}{}

{\def\LTcaptype{none} % do not increment counter
\begin{longtable}[]{@{}
  >{\raggedright\arraybackslash}p{(\linewidth - 4\tabcolsep) * \real{0.2108}}
  >{\raggedright\arraybackslash}p{(\linewidth - 4\tabcolsep) * \real{0.2108}}
  >{\raggedright\arraybackslash}p{(\linewidth - 4\tabcolsep) * \real{0.5785}}@{}}
\toprule\noalign{}
\begin{minipage}[b]{\linewidth}\raggedright
Дата
\end{minipage} & \begin{minipage}[b]{\linewidth}\raggedright
Версия
\end{minipage} & \begin{minipage}[b]{\linewidth}\raggedright
Содержание
\end{minipage} \\
\midrule\noalign{}
\endhead
\bottomrule\noalign{}
\endlastfoot
20.11.2026 & 1.0 & Подготовлен план по всей платформе. Включены карточки предметных областей, интерфейса, эксплуатации, сессий задач и планировщика \\
25.11.2026 & 1.1 & Самбьюи А.Ж.: уточнены TC-CFG-002, TC-CFG-006, TC-CFG-008 и TC-SYS-003: роли, представление просрочки, параметры публикации, повторные прохождения и восстановление. ID и остальные кейсы сохранены; исполнение не проводилось. Передача редактору DOC-10 и принятие ожидаются \\
\end{longtable}
}

\section{DOC-10. Верификация сессий задач и планировщика}\label{doc-10-ux432ux435ux440ux438ux444ux438ux43aux430ux446ux438ux44f-ux441ux435ux441ux441ux438ux439-ux437ux430ux434ux430ux447-ux438-ux43fux43bux430ux43dux438ux440ux43eux432ux449ux438ux43aux430}

\hypertarget{doc-10-верификация-сессий-задач-и-планировщика}{}

{\def\LTcaptype{none} % do not increment counter
\begin{longtable}[]{@{}
  >{\raggedright\arraybackslash}p{(\linewidth - 2\tabcolsep) * \real{0.3200}}
  >{\raggedright\arraybackslash}p{(\linewidth - 2\tabcolsep) * \real{0.6800}}@{}}
\toprule\noalign{}
\begin{minipage}[b]{\linewidth}\raggedright
Реквизит
\end{minipage} & \begin{minipage}[b]{\linewidth}\raggedright
Значение
\end{minipage} \\
\midrule\noalign{}
\endhead
\bottomrule\noalign{}
\endlastfoot
Документ & Авторский раздел плана верификации программного обеспечения \\
Версия и дата & 1.0 от 20.11.2026 \\
Автор раздела & Анохин Г.А. \\
Владелец сводного DOC-10 & Кокорев В.Д. \\
\end{longtable}
}

\subsection{Предмет и способ проверки}\label{ux43fux440ux435ux434ux43cux435ux442-ux438-ux441ux43fux43eux441ux43eux431-ux43fux440ux43eux432ux435ux440ux43aux438}

\hypertarget{предмет-и-способ-проверки}{}

Сессия --- личная работа студента с выданным набором задач. Каждая задача имеет своё условие и состояние. Планировщик --- часть программы, которая периодически находит задачи с истёкшим сроком и завершает их по правилам системы. Верификация в этой главе означает проверку описанного поведения по требованиям. Верификатор по шагам разбирает карточку с исходными данными, событиями и ожидаемыми результатами. Он устанавливает, целиком ли создан набор, к своей ли задаче отнесён ответ и сохранится ли итог при истечении срока или перезапуске.

Предметный проверяющий --- часть программы, которая проверяет решение задачи своего вида. Его вердикт может быть положительным или отрицательным. Когда вердикт зафиксирован, попытка решения завершена. Для разбора поведения сессии правильность учебных ответов задана заранее. Математические алгоритмы проверяют владельцы соответствующих подсистем. Правила составления карточек и чтения их номеров приведены в {\hyperlink{как-читать-обозначения-и-ответы-системы}{основной части плана}}.

Проверочная карточка имеет обозначение вида {\hyperlink{tc-ses-001}{TC-SES-001}}. Здесь TC означает проверку, SES --- сессии задач, а три цифры различают карточки внутри области. Читатель сопоставляет карточку с {\href{../04\_Системные\_требования/DOC-05\_Системные\_требования.pdf\#doc-05-требования-к-сессиям-задач}{системными требованиями к сессиям}}, {\href{../04\_Системные\_требования/DOC-05\_Системные\_требования.pdf\#doc-05-общесистемная-интеграция}{правилами интеграции}}, {\href{../07\_Программные\_требования/DOC-06\_Программные\_требования.pdf\#doc-06-сводные-программные-требования}{общими программными требованиями}} и {\href{../07\_Программные\_требования/DOC-06.2\_Сессии\_задач\_и\_планировщик.pdf\#doc-062-сессии-задач-и-планировщик}{требованиями к подсистеме сессий}}. Тело ответа --- данные, которые сервер возвращает клиенту вместе со статусом. Код ответа, состав тела и представление данных определяет {\href{../07\_Программные\_требования/DOC-06.2\_Сессии\_задач\_и\_планировщик.pdf\#doc-062-контракты-сессий-задач}{контракт операций}}.

Для каждого действия прослеживают ответ клиенту (программе, пославшей запрос), состояние задачи, значение её поля \texttt{result} с сохранённым вердиктом и состояние сессии. Клиент и сервер обмениваются запросами и ответами по протоколу HTTP. Если одинаковый внешний ответ допускает несколько внутренних исходов, верификатор прослеживает изменения в хранилище. Транзакция объединяет несколько изменений данных, которые сохраняются целиком либо отменяются вместе. Такая отмена называется откатом. Строка в таблице хранилища представляет один записанный объект. Блокировка записи согласует доступ нескольких обработчиков к одним данным, например когда два запроса пытаются создать набор одному студенту.

Для задачи \texttt{NEW} означает открытое задание, \texttt{SOLVED} --- верное решение, \texttt{FAILED} --- неверное решение или завершение по просрочке. Для сессии \texttt{NEW} означает, что остаются задачи без итога. Когда все задачи завершены, сессия получает \texttt{SOLVED}, если все решены верно, иначе --- \texttt{FAILED}. В JSON, текстовом формате передачи структурированных данных, \texttt{null} означает отсутствие значения. По полю \texttt{result} различают две причины состояния \texttt{FAILED}. При неверном решении есть отрицательный вердикт, при просрочке без вердикта значение равно \texttt{null}.

\subsubsection{Общие исходные данные}\label{ux43eux431ux449ux438ux435-ux438ux441ux445ux43eux434ux43dux44bux435-ux434ux430ux43dux43dux44bux435}

\hypertarget{общие-исходные-данные}{}

\begin{itemize}
\tightlist
\item
  А и Б --- условные учётные записи для разбора, обе с правами студента (роль \texttt{USER}). JWT --- подписанный токен, по которому сервер подтверждает личность пользователя. Каждый предъявляет свой токен в заголовке \texttt{Authorization:\ Bearer\ \textless{}токен\textgreater{}}. Отклонения указаны в карточках доступа. У А и Б номер учебной группы получен из подтверждённого университетского профиля. Для группы А включён «Практикум» --- именованный шаблон с двумя позициями для генерации задач, ранцевой и Эль-Гамаля. Шаблон задаёт виды и порядок задач, а условия формируют при создании набора. Если карточке нужен другой состав, он указан в её условиях.
\item
  Для каждой задачи даны её условие, один верный ответ и один предметно неверный ответ в допустимом формате. Варианты с изменённым полем, размером или участником разбирают отдельно, возвращаясь к исходному состоянию.
\item
  Запись операции содержит метод HTTP и путь. В карточках самостоятельной практики создать набор означает передать \texttt{POST /\allowbreak{}api/\allowbreak{}tasks/\allowbreak{}session}, указав имя шаблона в параметре \texttt{configuration}. UUID --- уникальный идентификатор объекта. Ответ отправляют через \texttt{PUT /\allowbreak{}api/\allowbreak{}tasks/\allowbreak{}submit} с UUID конкретной задачи в \texttt{statement\allowbreak{}Id}. Сессию читают через \texttt{GET /\allowbreak{}api/\allowbreak{}tasks/\allowbreak{}session/\allowbreak{}\{session\allowbreak{}Id\}}, заменяя \texttt{\{session\allowbreak{}Id\}} на UUID набора. Для текста JSON используется кодировка UTF-8. Допустимые поля и их типы перечислены в {\href{../07\_Программные\_требования/DOC-06.2\_Сессии\_задач\_и\_планировщик.pdf\#представление-сессии-и-задачи}{схеме данных сессии}}.
\item
  Серверное время создания обозначено \texttt{t0}. Время приводится в UTC (всемирном координированном времени), буква \texttt{Z} в конце обозначает нулевое смещение от UTC. Если карточка не задаёт другого времени, рассматривают самостоятельную практику. Для неё \texttt{t0 =\allowbreak{} 12:00:00.\allowbreak{}000Z}, общий срок \texttt{D =\allowbreak{} 12:20:00.\allowbreak{}000Z} того же дня. Для контрольного мероприятия решение заканчивается в более ранний из двух моментов --- через 20 минут после личного старта или при закрытии мероприятия. Отдельный пример приведён в {\hyperlink{tc-ses-017}{TC-SES-017}}. Время \texttt{t\_in} --- момент получения сервером \textbf{полного} запроса. Предел фиксации вердикта \texttt{L =\allowbreak{} t\_\allowbreak{}in + 60 секунд}. Время показано без даты для краткости. В ответах сервера проверяют полную дату, время UTC, миллисекунды и суффикс \texttt{Z}.
\item
  При успехе код 200 означает выполнение запроса, код 201 --- создание сессии. При отказе клиент получает числовой статус HTTP и машинный код причины. В теле ответа \texttt{code} обозначает причину, \texttt{message} поясняет её человеку, \texttt{timestamp} указывает время события, UUID в \texttt{requestId} позволяет найти конкретное обращение. В каждой карточке используют только указанные причины отказа. Их значения и соответствующие HTTP-статусы перечислены в {\href{../07\_Программные\_требования/DOC-06.2\_Сессии\_задач\_и\_планировщик.pdf\#коды-и-представление-ошибок}{таблице ошибок}}.
\end{itemize}

\subsection{Создание сессии}\label{ux441ux43eux437ux434ux430ux43dux438ux435-ux441ux435ux441ux441ux438ux438}

\hypertarget{создание-сессии}{}

\subsubsection{TC-SES-001. Создание и сохранение состава}\label{tc-ses-001-ux441ux43eux437ux434ux430ux43dux438ux435-ux438-ux441ux43eux445ux440ux430ux43dux435ux43dux438ux435-ux441ux43eux441ux442ux430ux432ux430}

\hypertarget{tc-ses-001-создание-и-сохранение-состава}{}

\hypertarget{tc-ses-001}{}

\textbf{Исходные условия.} У А нет активной сессии. Для этой карточки «Практикум» содержит три позиции. Первая и третья --- ранцевые задачи с разными условиями, вторая --- задача Эль-Гамаля. Известны условия, которые вернёт генератор для каждой позиции.

\textbf{Действия.} В момент \texttt{t0} создать «Практикум». Сопоставить тип и условие каждой задачи с её позицией в шаблоне. Записать UUID в порядке выдачи. Изменить состав и порядок задач шаблона группы, перезапустить приложение и до срока прочитать выданную сессию.

\textbf{Ожидаемый результат.} Ответ на создание имеет статус 201. Владелец сессии --- А, время создания равно \texttt{t0}. Поле \texttt{assessment\allowbreak{}Id} отсылает к контрольному мероприятию и здесь равно \texttt{null}. Участие в контрольном мероприятии и его оценка не появляются. Сессия и все три задачи имеют состояние \texttt{NEW}, у каждой задачи \texttt{result\ =\ null}. В списке сохранён исходный порядок типов. Первая и третья позиции содержат отдельные задачи, каждая со своим UUID и своим условием. У всех задач \texttt{sessionId} совпадает с \texttt{id} сессии. Повторное чтение даёт 200 и возвращает ту же последовательность UUID. На каждой позиции сохранены её тип и условие. Поле \texttt{statement} имеет строковый тип. После разбора этой строки как JSON получают исходное условие задачи. Владелец, время создания и срок \texttt{D} сохраняются.

\subsubsection{TC-SES-002. Доступ пользователя к шаблону}\label{tc-ses-002-ux434ux43eux441ux442ux443ux43f-ux43fux43eux43bux44cux437ux43eux432ux430ux442ux435ux43bux44f-ux43a-ux448ux430ux431ux43bux43eux43dux443}

\hypertarget{tc-ses-002-доступ-пользователя-к-шаблону}{}

\hypertarget{tc-ses-002}{}

\textbf{Исходные условия.} Активной сессии нет. Каждый вариант запроса начинается с исходного состояния.

\textbf{Действия и ожидаемые результаты.} Аутентификация подтверждает, от какого лица поступил запрос. Запрос создания без действительной аутентификации даёт 401 \texttt{UNAUTHENTICATED}, с действительной аутентификацией без роли \texttt{USER} --- 403 \texttt{ACCESS\_\allowbreak{}DENIED}. Для пользователя без подтверждённой группы, отключённого шаблона, имени шаблона другой группы и имени с другим регистром ожидается 404 \texttt{CONFIGURATION\_\allowbreak{}NOT\_\allowbreak{}FOUND}. В каждом случае набор отсутствует. При верных полномочиях и имени новый набор принадлежит пользователю, чья личность подтверждена при входе.

\subsubsection{TC-SES-003. Конкурирующие создания и потеря ответа}\label{tc-ses-003-ux43aux43eux43dux43aux443ux440ux438ux440ux443ux44eux449ux438ux435-ux441ux43eux437ux434ux430ux43dux438ux44f-ux438-ux43fux43eux442ux435ux440ux44f-ux43eux442ux432ux435ux442ux430}

\hypertarget{tc-ses-003-конкурирующие-создания-и-потеря-ответа}{}

\hypertarget{tc-ses-003}{}

\textbf{Исходные условия.} У А нет сессии \texttt{NEW}. Обе копии запроса содержат одно и то же допустимое имя шаблона.

\textbf{Действия.} Разобрать два порядка событий. В первом ответ на создание теряется по пути к клиенту. А повторяет запрос до \texttt{D}, пока созданная сессия имеет состояние \texttt{NEW}. Во втором оба запроса поступают до того, как первый завершил создание.

\textbf{Ожидаемый результат.} Существует одна новая сессия. Один запрос создаёт её со статусом 201, другой получает 409 \texttt{ACTIVE\_\allowbreak{}SESSION\_\allowbreak{}EXISTS}. В \texttt{details.\allowbreak{}session\allowbreak{}Id} указан UUID единственного созданного набора. Чтение с этим UUID возвращает весь набор пользователю А. При одновременных запросах на создание проверка наличия сессии и вставка нового набора происходят под блокировкой строки одного пользователя в одной транзакции.

\subsubsection{TC-SES-004. Сбой при создании полного набора}\label{tc-ses-004-ux441ux431ux43eux439-ux43fux440ux438-ux441ux43eux437ux434ux430ux43dux438ux438-ux43fux43eux43bux43dux43eux433ux43e-ux43dux430ux431ux43eux440ux430}

\hypertarget{tc-ses-004-сбой-при-создании-полного-набора}{}

\hypertarget{tc-ses-004}{}

\textbf{Исходные условия.} Первое условие уже сформировано. Во втором элементе шаблона задан неподдерживаемый тип. В отдельном варианте оба условия сформированы, но запись набора в хранилище завершается ошибкой.

\textbf{Действия.} Запросить создание в каждом варианте. Исправить шаблон либо восстановить запись, повторить исходный запрос.

\textbf{Ожидаемый результат.} Первая операция даёт 500 \texttt{INTERNAL\_\allowbreak{}ERROR}. Данные незавершённого набора и его задач отсутствуют, ответ 201 не выдан. После устранения причины повтор возвращает 201 с обеими задачами. По описанию транзакции создание записей сессии и задач либо завершается целиком, либо целиком откатывается.

\subsection{Приём решения и предметный результат}\label{ux43fux440ux438ux451ux43c-ux440ux435ux448ux435ux43dux438ux44f-ux438-ux43fux440ux435ux434ux43cux435ux442ux43dux44bux439-ux440ux435ux437ux443ux43bux44cux442ux430ux442}

\hypertarget{приём-решения-и-предметный-результат}{}

\subsubsection{TC-SES-005. Ошибка формы запроса и граница размера}\label{tc-ses-005-ux43eux448ux438ux431ux43aux430-ux444ux43eux440ux43cux44b-ux437ux430ux43fux440ux43eux441ux430-ux438-ux433ux440ux430ux43dux438ux446ux430-ux440ux430ux437ux43cux435ux440ux430}

\hypertarget{tc-ses-005-ошибка-формы-запроса-и-граница-размера}{}

\hypertarget{tc-ses-005}{}

\textbf{Исходные условия.} А владеет задачей Эль-Гамаля в состоянии \texttt{NEW}. Схема предметного ответа описывает, какие поля обязательны и какого типа должно быть каждое. Поля \texttt{k} для Эль-Гамаля (одноразовый показатель) и \texttt{power} для ранцевой задачи (размерность) приведены в {\href{../07\_Программные\_требования/DOC-06.2\_Сессии\_задач\_и\_планировщик.pdf\#предметный-ответ}{схемах решений}}. До \texttt{D} достаточно времени для исправления ответа.

\textbf{Действия и ожидаемые результаты.}

\begin{enumerate}
\def\labelenumi{\arabic{enumi}.}
\tightlist
\item
  При создании у другого пользователя с доступным шаблоном повторить параметр \texttt{configuration}, передать пустое имя, передать одновременно \texttt{configuration} и \texttt{assessment\allowbreak{}Id} либо опустить оба параметра. При отправке решения А пропустить \texttt{statement\allowbreak{}Id}, повторить его или передать строку, не являющуюся UUID. В каждом варианте ожидается 400 \texttt{INVALID\_\allowbreak{}REQUEST}.
\item
  Для Эль-Гамаля передать пустое тело, ошибочный JSON, тело без \texttt{k}, с \texttt{k\ =\ null}, числом вместо десятичной строки или строкой с ведущим нулём. Для ранцевой задачи проверить вариант с обязательным полем \texttt{power} строкового типа. Каждый из этих запросов даёт 400 \texttt{INVALID\_\allowbreak{}REQUEST}. Вызов предметного проверяющего отсутствует.
\item
  Дополнить заведомо верный JSON пробелами после закрывающей скобки до размера 1 048 577 байт. Ожидается 413 \texttt{PAYLOAD\_\allowbreak{}TOO\_\allowbreak{}LARGE}. После всех отказов попытка доступна, состояние задачи \texttt{NEW}, результат равен \texttt{null}.
\item
  Отправить тот же верный JSON с общим размером ровно 1 048 576 байт. Запрос укладывается в предел, проходит предметную проверку и возвращает 200 с \texttt{is\allowbreak{}Ok =\allowbreak{} true}.
\end{enumerate}

\subsubsection{TC-SES-006. Выбор проверяющего и итог нескольких задач}\label{tc-ses-006-ux432ux44bux431ux43eux440-ux43fux440ux43eux432ux435ux440ux44fux44eux449ux435ux433ux43e-ux438-ux438ux442ux43eux433-ux43dux435ux441ux43aux43eux43bux44cux43aux438ux445-ux437ux430ux434ux430ux447}

\hypertarget{tc-ses-006-выбор-проверяющего-и-итог-нескольких-задач}{}

\hypertarget{tc-ses-006}{}

\textbf{Исходные условия.} А получил набор с двумя задачами. Выданные условия и UUID записаны. Шаблон группы после выдачи изменён. До срока оба ответа успевают получить вердикт.

\textbf{Действия.} Отправить допустимые по форме решения для исходных условий. Для первой задачи получить отрицательный вердикт, затем для второй --- положительный. Разобрать также одновременное сохранение вердиктов разных задач с каждым из двух возможных порядков фиксации. На отдельном наборе обе задачи решить верно.

\textbf{Ожидаемый результат.} Обработчик каждого типа получает своё сохранённое условие и переданное решение. Ответ на каждый вердикт имеет статус 200. Поля \texttt{problemId} и \texttt{problem\allowbreak{}Type} соответствуют выданной задаче. При ошибочном решении \texttt{is\allowbreak{}Ok =\allowbreak{} false}, \texttt{error\allowbreak{}Message} содержит пояснение, задача переходит в \texttt{FAILED}. При верном --- \texttt{is\allowbreak{}Ok =\allowbreak{} true}, \texttt{error\allowbreak{}Message =\allowbreak{} null}, состояние \texttt{SOLVED}. После первого сохранения сессия ещё \texttt{NEW}. После обоих сохранившийся итог смешанного набора равен \texttt{FAILED}, а набора с двумя верными решениями --- \texttt{SOLVED}. Результат одновременной записи совпадает с результатом последовательной для тех же двух вердиктов.

\textbf{Вариант с пользовательским заданием.} Проверить две отдельные выдачи позиции \texttt{CUSTOM\_\allowbreak{}TEXT}. Для разбираемого варианта текст условия --- «Назовите столицу России», сохранённый эталон --- \texttt{Москва}. При чтении условия эталон отсутствует. В первой выдаче отправить \texttt{\{"answer":"\ \ Москва\ \ "\}}. Ожидается \texttt{200} с \texttt{is\allowbreak{}Ok =\allowbreak{} true}. Во второй выдаче того же варианта отправить \texttt{\{"answer":"москва"\}}. Ожидается \texttt{200} с \texttt{is\allowbreak{}Ok =\allowbreak{} false}. В обоих случаях сохраняется единственный вердикт по UUID задачи соответствующего набора.

\subsubsection{TC-SES-007. Две посылки на одну задачу}\label{tc-ses-007-ux434ux432ux435-ux43fux43eux441ux44bux43bux43aux438-ux43dux430-ux43eux434ux43dux443-ux437ux430ux434ux430ux447ux443}

\hypertarget{tc-ses-007-две-посылки-на-одну-задачу}{}

\hypertarget{tc-ses-007}{}

\textbf{Исходные условия.} Один допустимый ответ А принят до \texttt{D}. Его запись ещё открыта. Верный ответ и второй, неверный, имеют разные значения.

\textbf{Действия.} До \texttt{D} передать второй ответ, пока обрабатывается первый. Завершить первый положительным вердиктом. Ещё раз до \texttt{D} отправить второй ответ и прочитать набор.

\textbf{Ожидаемый результат.} Повтор при открытой операции получает 409 \texttt{CHECK\_\allowbreak{}IN\_\allowbreak{}PROGRESS}, после сохранения вердикта --- 409 \texttt{TASK\_\allowbreak{}FINALIZED}. Итог получен для исходного ответа ровно один раз. Сохранённые тело первой операции и положительный результат остаются прежними.

\textbf{Одновременное поступление.} В отдельном наборе зарегистрировать два запроса владельца на одну свободную задачу до \texttt{D}. Оба тела имеют допустимый формат. Одно решение верно, другое неверно. Подвести оба запроса к началу предметной проверки одновременно. Проверка доступности задачи и закрепление за ней одной операции выполняются как одно неделимое действие. Право начать вычисление получает только эта операция. Пока она выполняется, вторая получает 409 \texttt{CHECK\_\allowbreak{}IN\_\allowbreak{}PROGRESS}. Если право получило верное решение, задача становится \texttt{SOLVED}, если неверное --- \texttt{FAILED} с отрицательным результатом. В обоих порядках проверяющему передано только тело выбранной операции. Последующее чтение возвращает единственный соответствующий ему вердикт.

\textbf{Отказ одному из двух запросов.} Два ответа владельца поступили до \texttt{D}. После \texttt{D} при разборе первого обнаружилась ошибка формата. Он получает 400 \texttt{INVALID\_\allowbreak{}REQUEST}. Пока второй своевременный ответ ждёт проверки и его собственный предел ещё не достигнут, задача остаётся \texttt{NEW}. До этого предела вердикт по второму ответу может быть зафиксирован.

\subsubsection{TC-SES-008. Принадлежность объекта и порядок отказов}\label{tc-ses-008-ux43fux440ux438ux43dux430ux434ux43bux435ux436ux43dux43eux441ux442ux44c-ux43eux431ux44aux435ux43aux442ux430-ux438-ux43fux43eux440ux44fux434ux43eux43a-ux43eux442ux43aux430ux437ux43eux432}

\hypertarget{tc-ses-008-принадлежность-объекта-и-порядок-отказов}{}

\hypertarget{tc-ses-008}{}

\textbf{Исходные условия.} А владеет известной задачей. У Б есть собственная учётная запись с ролью \texttt{USER}. Для сравнения взяты UUID отсутствующих задачи и сессии.

\textbf{Действия.} Б отправляет решение на задачу А и на отсутствующую задачу. Б читает сессию А и отсутствующую сессию. Повторить обе операции без действительных учётных данных и с учётной записью без роли \texttt{USER}. В отдельном варианте Б обращается к задаче А точно в \texttt{D}.

\textbf{Ожидаемый результат.} Для Б оба обращения к задачам дают 404 \texttt{TASK\_\allowbreak{}NOT\_\allowbreak{}FOUND}, оба чтения --- 404 \texttt{SESSION\_\allowbreak{}NOT\_\allowbreak{}FOUND}. Тот же 404 сохраняется и для чужой задачи в момент \texttt{D}. Для обеих операций отсутствие аутентификации даёт 401 \texttt{UNAUTHENTICATED}, отсутствие роли --- 403 \texttt{ACCESS\_\allowbreak{}DENIED}. Проверяющий не вызывается. До срока состояние и результаты сессии А сохраняются, после срока для них действует обычный порядок завершения.

\textbf{Просрочка после чужой отправки.} В отдельном наборе у А одна задача \texttt{NEW}, свой ответ А не посылал. В 12:19:59.999Z Б отправляет ответ на эту задачу. Сравнить с отправкой Б на отсутствующий UUID в тот же момент. В обоих случаях ожидается 404 \texttt{TASK\_\allowbreak{}NOT\_\allowbreak{}FOUND}. В отдельном варианте остановить приложение после получения полного запроса Б, до сверки принадлежности и открытия операции. При обрыве соединения исходный HTTP-ответ неизвестен. После перезапуска за границей \texttt{D} открытой операции Б по задаче А нет. Повторное обращение Б даёт 404 \texttt{TASK\_\allowbreak{}NOT\_\allowbreak{}FOUND}. Задача А завершается по сроку. После \texttt{D} при обращении А за новым набором прежняя задача и её сессия завершаются в \texttt{FAILED} с \texttt{result\ =\ null} у задачи. Создание возвращает 201 с новым UUID сессии.

\subsection{Срок сессии и срок обработки ответа}\label{ux441ux440ux43eux43a-ux441ux435ux441ux441ux438ux438-ux438-ux441ux440ux43eux43a-ux43eux431ux440ux430ux431ux43eux442ux43aux438-ux43eux442ux432ux435ux442ux430}

\hypertarget{срок-сессии-и-срок-обработки-ответа}{}

\subsubsection{TC-SES-009. Своевременный ответ в очереди во время просрочки}\label{tc-ses-009-ux441ux432ux43eux435ux432ux440ux435ux43cux435ux43dux43dux44bux439-ux43eux442ux432ux435ux442-ux432-ux43eux447ux435ux440ux435ux434ux438-ux432ux43e-ux432ux440ux435ux43cux44f-ux43fux440ux43eux441ux440ux43eux447ux43aux438}

\hypertarget{tc-ses-009-своевременный-ответ-в-очереди-во-время-просрочки}{}

\hypertarget{tc-ses-009}{}

\textbf{Исходные условия.} У А две задачи \texttt{NEW}. Запрос на первую задачу ожидает в очереди, то есть уже зарегистрирован и ждёт освобождения проверяющего. Полный запрос получен в 12:19:59.900Z. Предметная проверка начнётся в 12:20:02.000Z и зафиксирует верный результат в 12:20:04.000Z. По второй задаче ответа нет.

\textbf{Действия.} В момент \texttt{D} рассмотреть запуск планировщика между получением запроса и запуском проверяющего. До окончания проверки прочитать сессию и запросить новый набор. После фиксации вердикта снова прочитать сессию.

\textbf{Ожидаемый результат.} Планировщик сохраняет первую задачу в \texttt{NEW}, вторую завершает по сроку в \texttt{FAILED} с \texttt{result\ =\ null}. Чтение видит сессию \texttt{NEW}, создание получает 409 \texttt{ACTIVE\_\allowbreak{}SESSION\_\allowbreak{}EXISTS}. Фиксация вердикта первой задачи происходит до её предела 12:20:59.900Z, поэтому она становится \texttt{SOLVED}, итог сессии --- \texttt{FAILED}. Время запуска проверяющего и фоновой обработки не меняет \texttt{t\_in}.

\subsubsection{TC-SES-010. Полный запрос на границе срока}\label{tc-ses-010-ux43fux43eux43bux43dux44bux439-ux437ux430ux43fux440ux43eux441-ux43dux430-ux433ux440ux430ux43dux438ux446ux435-ux441ux440ux43eux43aux430}

\hypertarget{tc-ses-010-полный-запрос-на-границе-срока}{}

\hypertarget{tc-ses-010}{}

\textbf{Исходные условия.} Одна собственная задача уже имеет сохранённый вердикт, другая остаётся \texttt{NEW}. Ближайший плановый проход назначен на 12:20:30.000Z. Для каждого варианта берут исходное состояние.

\textbf{Действия.} Начать передачу решения за одну миллисекунду до \texttt{D}, передать последний байт в \texttt{D}. Повторить вариант с полным запросом в \texttt{D\ +\ 1\ миллисекунда}. Отдельно направить такой запрос к уже завершённой собственной задаче. До фонового прохода прочитать набор. В другом варианте сразу после \texttt{D} первым запросом создать новую сессию.

\textbf{Ожидаемый результат.} Каждая поздняя отправка получает 410 \texttt{SESSION\_\allowbreak{}EXPIRED}, в том числе для завершённой собственной задачи. Отметка клиента и время начала передачи не влияют на допуск. Чтение возвращает прежний вердикт и просроченную задачу в \texttt{FAILED} с \texttt{result\ =\ null}. Итог набора \texttt{FAILED}. Новый набор может быть создан после завершения прежнего, если тот уже просрочен и своевременных обработок нет. Повторный фоновый проход не меняет прежних итогов. По расписанию подсистемы плановый период проходов равен 60 секундам.

\subsubsection{TC-SES-011. Строгая граница фиксации вердикта}\label{tc-ses-011-ux441ux442ux440ux43eux433ux430ux44f-ux433ux440ux430ux43dux438ux446ux430-ux444ux438ux43aux441ux430ux446ux438ux438-ux432ux435ux440ux434ux438ux43aux442ux430}

\hypertarget{tc-ses-011-строгая-граница-фиксации-вердикта}{}

\hypertarget{tc-ses-011}{}

\textbf{Исходные условия.} Для каждого варианта у А сессия с одной задачей в состоянии \texttt{NEW}. Верный ответ принят в 12:05:00.000Z, поэтому \texttt{L =\allowbreak{} 12:06:00.\allowbreak{}000Z}. Очередь и расчёт занимают часть этого интервала.

\textbf{Действия и ожидаемые результаты.} При фиксации в 12:05:59.999Z верный вердикт сохраняется, задача и сессия получают \texttt{SOLVED}. Если вычисление закончено раньше, но к 12:06:00.000Z вердикт ещё не записан, операция завершается с 503 \texttt{CHECK\_\allowbreak{}TIMEOUT}. То же ожидается при поступлении результата в 12:06:00.001Z. Поздний результат отбрасывается. Задача остаётся \texttt{NEW} с \texttt{result\ =\ null}, поскольку общий срок ещё не наступил. Новый правильный ответ до \texttt{D} может получить 200 и завершить задачу.

\subsubsection{TC-SES-012. Отказ до и после общего срока}\label{tc-ses-012-ux43eux442ux43aux430ux437-ux434ux43e-ux438-ux43fux43eux441ux43bux435-ux43eux431ux449ux435ux433ux43e-ux441ux440ux43eux43aux430}

\hypertarget{tc-ses-012-отказ-до-и-после-общего-срока}{}

\hypertarget{tc-ses-012}{}

\textbf{Исходные условия.} У А одна задача. Первый вариант начинается в 12:05:00.000Z. Во втором полный запрос получен в 12:19:59.500Z. Принимающая сторона обнаруживает ошибку формата либо получает отказ проверяющего без вердикта.

\textbf{Действия.} Для раннего запроса получить отказ через секунду. Исправить формат либо восстановить доступность проверяющего и послать новый верный ответ до \texttt{D}. Для позднего варианта обнаружить отказ в 12:20:00.500Z, затем прочитать задачу и повторить отправку.

\textbf{Ожидаемый результат.} Ошибка формата даёт 400 \texttt{INVALID\_\allowbreak{}REQUEST}, недоступность проверяющего --- 503 \texttt{CHECK\_\allowbreak{}UNAVAILABLE}. В первом варианте после отказа видны \texttt{NEW} и \texttt{result\ =\ null}. Исправленный ответ может завершить задачу. Во втором ни один отказ не образует предметного результата. Задача и сессия завершаются по просрочке в \texttt{FAILED}, чтение возвращает \texttt{result\ =\ null}, новый ответ получает 410 \texttt{SESSION\_\allowbreak{}EXPIRED}.

\subsection{Восстановление и согласованность записи}\label{ux432ux43eux441ux441ux442ux430ux43dux43eux432ux43bux435ux43dux438ux435-ux438-ux441ux43eux433ux43bux430ux441ux43eux432ux430ux43dux43dux43eux441ux442ux44c-ux437ux430ux43fux438ux441ux438}

\hypertarget{восстановление-и-согласованность-записи}{}

\subsubsection{TC-SES-013. Перезапуск с незавершённой операцией}\label{tc-ses-013-ux43fux435ux440ux435ux437ux430ux43fux443ux441ux43a-ux441-ux43dux435ux437ux430ux432ux435ux440ux448ux451ux43dux43dux43eux439-ux43eux43fux435ux440ux430ux446ux438ux435ux439}

\hypertarget{tc-ses-013-перезапуск-с-незавершённой-операцией}{}

\hypertarget{tc-ses-013}{}

\textbf{Исходные условия.} В сессии А одна задача Эль-Гамаля. Назначения полей и исходное условие приведены в {\href{../07\_Программные\_требования/DOC-06.2\_Сессии\_задач\_и\_планировщик.pdf\#учебный-пример-решения-эль-гамаля}{учебном примере Эль-Гамаля}}. Для этого условия известен верный ответ с полями \texttt{y\ =\ "8"}, \texttt{k\ =\ "3"}, \texttt{c1\ =\ "10"}, \texttt{c2\ =\ "9"}, \texttt{decrypted\allowbreak{}Message =\allowbreak{} "13"}. Время получения полного запроса \texttt{t\_\allowbreak{}in =\allowbreak{} 12:19:59.\allowbreak{}000Z}, предел \texttt{L =\allowbreak{} 12:20:59.\allowbreak{}000Z}. Открытая запись операции уже сохранена. Приложение останавливается, пока ответ ожидает вызова предметного проверяющего.

\textbf{Действия.} В первом варианте восстановить работу в 12:20:30.000Z. Два обработчика восстановления одновременно выбирают ту же открытую запись. Завершить верную проверку в 12:20:58.999Z. Во втором варианте восстановить приложение ровно в 12:20:59.000Z. Для каждого варианта прочитать набор.

\textbf{Ожидаемый результат.} Из записи восстанавливаются UUID операции, владелец А, задача, все пять полей ответа с исходными значениями, время получения 12:19:59.000Z и предел 12:20:59.000Z. В первом варианте запись для обработки получает один из двух обработчиков. Он передаёт сохранённое условие и восстановленное тело проверяющему, сохраняет один положительный вердикт. Чтение возвращает задачу и сессию в состоянии \texttt{SOLVED}. Во втором варианте запись получает технический исход \texttt{CHECK\_\allowbreak{}TIMEOUT}, даже если очередь после запуска свободна. Состояние задачи и сессии --- \texttt{FAILED} по просрочке, \texttt{result\ =\ null}. Перезапуск не продлевает ни один из сроков.

\subsubsection{TC-SES-014. Потеря ответа после сохранения вердикта}\label{tc-ses-014-ux43fux43eux442ux435ux440ux44f-ux43eux442ux432ux435ux442ux430-ux43fux43eux441ux43bux435-ux441ux43eux445ux440ux430ux43dux435ux43dux438ux44f-ux432ux435ux440ux434ux438ux43aux442ux430}

\hypertarget{tc-ses-014-потеря-ответа-после-сохранения-вердикта}{}

\hypertarget{tc-ses-014}{}

\textbf{Исходные условия.} У А две задачи. Состояние первой --- \texttt{SOLVED}. Для второй до срока сохранён предметно неверный ответ, но соединение закрылось до доставки 200.

\textbf{Действия.} Перезапустить приложение после фиксации второго вердикта. Считать собственную сессию дважды, в том числе после \texttt{D}. Сопоставить содержимое с выданными при создании идентификаторами, условиями и временем.

\textbf{Ожидаемый результат.} Оба чтения возвращают 200. Сохранены владелец, \texttt{createdAt}, набор условий, типы, UUID и связи задач. У первой задачи \texttt{SOLVED} и положительный результат, у второй \texttt{FAILED} и отрицательный результат с её UUID, типом и непустым пояснением. Сессия \texttt{FAILED}. Ранее закрытая операция не запускается повторно. Чтение не расходует попытку и возвращает один и тот же предметный результат.

\subsubsection{TC-SES-015. Сбой при фиксации вердикта}\label{tc-ses-015-ux441ux431ux43eux439-ux43fux440ux438-ux444ux438ux43aux441ux430ux446ux438ux438-ux432ux435ux440ux434ux438ux43aux442ux430}

\hypertarget{tc-ses-015-сбой-при-фиксации-вердикта}{}

\hypertarget{tc-ses-015}{}

\textbf{Исходные условия.} У А сессия с одной задачей \texttt{NEW}. До \texttt{D} и \texttt{L} проверяющий передал положительный вердикт. При записи произошёл сбой. Сравнивают два независимых исхода.

\textbf{Отказ сообщён клиенту.} После отката фиксации операция закрыта с техническим отказом, клиент получил 500 \texttt{INTERNAL\_\allowbreak{}ERROR}. Прочитать сессию и отправить новый ответ до \texttt{D}. До новой отправки задача и сессия остаются \texttt{NEW}, \texttt{result\ =\ null}. Первоначальная операция не возобновляется, и её вердикт не может появиться позже. Для новой отправки доступна предметная попытка.

\textbf{Работа прервана до выдачи ответа.} После отката в хранилище остаётся открытая запись принятого запроса, клиент не получил ни 200, ни сообщения об отказе. Восстановить приложение до исходного \texttt{L}. До восстановления задача и сессия имеют состояние \texttt{NEW}, \texttt{result\ =\ null}. Обработка продолжается по данным той же записи. Если вердикт зафиксирован до \texttt{L}, чтение возвращает один положительный результат и состояние \texttt{SOLVED} у задачи и сессии. При достижении \texttt{L} без фиксации применяется \texttt{CHECK\_\allowbreak{}TIMEOUT}. Вердикт, оба состояния и закрытие успешной операции записываются вместе.

\subsubsection{TC-SES-016. Граница транзакции при одновременном завершении}\label{tc-ses-016-ux433ux440ux430ux43dux438ux446ux430-ux442ux440ux430ux43dux437ux430ux43aux446ux438ux438-ux43fux440ux438-ux43eux434ux43dux43eux432ux440ux435ux43cux435ux43dux43dux43eux43c-ux437ux430ux432ux435ux440ux448ux435ux43dux438ux438}

\hypertarget{tc-ses-016-граница-транзакции-при-одновременном-завершении}{}

\hypertarget{tc-ses-016}{}

\textbf{Исходные условия.} Две последние задачи одной сессии получили верный и неверный предметные вердикты. Оба обработчика готовы их фиксировать. Варианты внешнего результата описаны в {\hyperlink{tc-ses-006}{TC-SES-006}}.

\textbf{Действия.} По описанию подсистемы проследить, какие строки блокирует каждый обработчик во время предметного вычисления и при записи. Разобрать смену очередности двух записей и одновременное обращение планировщика к той же задаче с уже принятым своевременным ответом.

\textbf{Ожидаемый результат.} Предметное вычисление вынесено за пределы транзакции фиксации. При записи сначала блокируется строка сессии, затем задачи, затем операции. После получения блокировки обработчик учитывает актуальные состояния. Вердикт, состояния задачи и сессии и закрытие операции фиксируются одной транзакцией. Ни одна очередность не оставляет сессию \texttt{NEW} после завершения последней задачи и не позволяет планировщику заместить своевременный вердикт. В обоих вариантах сохраняются оба результата, итог сессии --- \texttt{FAILED}.

\subsubsection{TC-SES-017. Общий срок личной сессии внутри контрольного окна}\label{tc-ses-017-ux43eux431ux449ux438ux439-ux441ux440ux43eux43a-ux43bux438ux447ux43dux43eux439-ux441ux435ux441ux441ux438ux438-ux432ux43dux443ux442ux440ux438-ux43aux43eux43dux442ux440ux43eux43bux44cux43dux43eux433ux43e-ux43eux43aux43dux430}

\hypertarget{tc-ses-017-общий-срок-личной-сессии-внутри-контрольного-окна}{}

\hypertarget{tc-ses-017}{}

\textbf{Условия.} Мероприятие группы опубликовано с \texttt{opens\allowbreak{}At =\allowbreak{} 10:00:00.\allowbreak{}000Z}, \texttt{closes\allowbreak{}At =\allowbreak{} 11:00:00.\allowbreak{}000Z}. Личное время старта обозначено \texttt{t0}, срок решения --- \texttt{D}, фиксация вердикта должна произойти строго до предела \texttt{L =\allowbreak{} t\_\allowbreak{}in + 60 секунд}. Каждый вариант начинают на независимых данных.

\textbf{Действия.} Один студент начинает контрольную сессию в 10:05, другой --- в 10:55. Для сравнения третий студент начинает самостоятельную практику в 10:55. Прочитать у каждого \texttt{deadlineAt} --- сохранённое время окончания решения. Во второй контрольной сессии до самого закрытия остаются две задачи. На одну поступает полный ответ в 10:59:59.900Z, на другую --- ровно в 11:00:00.000Z. Провести контроль просрочки в 11:00 и рассмотреть положительный вердикт по первому ответу при фиксации в 11:00:02.000Z. Отдельно разобрать отсутствие фиксации именно к исходному пределу 11:00:59.900Z.

\textbf{Ожидаемый результат.} Для контроля \texttt{D =\allowbreak{} min(t0 + 20 минут, closes\allowbreak{}At)}. В первой контрольной сессии срок равен 10:25. Во второй срок равен 11:00. У начатой в 10:55 практики \texttt{D\ =\ 11:15}. Повторное чтение, в том числе после перезапуска, возвращает прежние значения. Своевременно зарегистрированный ответ второй контрольной сессии защищён от фоновой просрочки до своего прежнего \texttt{L}. При фиксации в 11:00:02 задача получает \texttt{SOLVED} со своим вердиктом. Новый ответ на другую задачу, полученный ровно в \texttt{D} или позже, получает 410 \texttt{SESSION\_\allowbreak{}EXPIRED}. При отсутствии другого своевременного ответа задача завершается в \texttt{FAILED} с \texttt{result\ =\ null}. Если первый вердикт также не зафиксирован до \texttt{L}, операция завершается с 503 \texttt{CHECK\_\allowbreak{}TIMEOUT}, а задача --- по просрочке без предметного результата. Закрытие мероприятия ограничивает приём новых ответов. Оно не сокращает и не продлевает срок фиксации уже зарегистрированного допустимого ответа.

\subsection{Соответствие карточек требованиям}\label{ux441ux43eux43eux442ux432ux435ux442ux441ux442ux432ux438ux435-ux43aux430ux440ux442ux43eux447ux435ux43a-ux442ux440ux435ux431ux43eux432ux430ux43dux438ux44fux43c}

\hypertarget{соответствие-карточек-требованиям}{}

Таблица даёт владельцу матрицы DOC-09 связи по домену сессий. Каждая строка обозначает карточку и требования, для которых в ней указан наблюдаемый результат или признак при разборе проекта. Системные обязательства по интеграции представлены в пределах операций сессий.

{\def\LTcaptype{none} % do not increment counter
\begin{longtable}[]{@{}
  >{\raggedright\arraybackslash}p{(\linewidth - 4\tabcolsep) * \real{0.1993}}
  >{\raggedright\arraybackslash}p{(\linewidth - 4\tabcolsep) * \real{0.3734}}
  >{\raggedright\arraybackslash}p{(\linewidth - 4\tabcolsep) * \real{0.4273}}@{}}
\toprule\noalign{}
\begin{minipage}[b]{\linewidth}\raggedright
Карточка
\end{minipage} & \begin{minipage}[b]{\linewidth}\raggedright
Системные требования
\end{minipage} & \begin{minipage}[b]{\linewidth}\raggedright
Программные требования
\end{minipage} \\
\midrule\noalign{}
\endhead
\bottomrule\noalign{}
\endlastfoot
{\hyperlink{tc-ses-001}{TC-SES-001}} & {\href{../04\_Системные\_требования/DOC-05\_Системные\_требования.pdf\#sr-ses-001}{SR-SES-001}}, {\href{../04\_Системные\_требования/DOC-05\_Системные\_требования.pdf\#sr-ses-006}{SR-SES-006}}, {\href{../04\_Системные\_требования/DOC-05\_Системные\_требования.pdf\#sr-sys-002}{SR-SYS-002}}, {\href{../04\_Системные\_требования/DOC-05\_Системные\_требования.pdf\#sr-sys-003}{SR-SYS-003}}, {\href{../04\_Системные\_требования/DOC-05\_Системные\_требования.pdf\#sr-sys-007}{SR-SYS-007}} & {\href{../07\_Программные\_требования/DOC-06.2\_Сессии\_задач\_и\_планировщик.pdf\#swr-ses-002}{SWR-SES-002}}, {\href{../07\_Программные\_требования/DOC-06.2\_Сессии\_задач\_и\_планировщик.pdf\#swr-ses-004}{SWR-SES-004}}, {\href{../07\_Программные\_требования/DOC-06.2\_Сессии\_задач\_и\_планировщик.pdf\#swr-ses-006}{SWR-SES-006}}, {\href{../07\_Программные\_требования/DOC-06.2\_Сессии\_задач\_и\_планировщик.pdf\#swr-ses-019}{SWR-SES-019}}, {\href{../07\_Программные\_требования/DOC-06\_Программные\_требования.pdf\#swr-sys-002}{SWR-SYS-002}}, {\href{../07\_Программные\_требования/DOC-06\_Программные\_требования.pdf\#swr-sys-003}{SWR-SYS-003}}, {\href{../07\_Программные\_требования/DOC-06\_Программные\_требования.pdf\#swr-sys-008}{SWR-SYS-008}} \\
{\hyperlink{tc-ses-002}{TC-SES-002}} & {\href{../04\_Системные\_требования/DOC-05\_Системные\_требования.pdf\#sr-ses-001}{SR-SES-001}}, {\href{../04\_Системные\_требования/DOC-05\_Системные\_требования.pdf\#sr-sys-001}{SR-SYS-001}}, {\href{../04\_Системные\_требования/DOC-05\_Системные\_требования.pdf\#sr-sys-003}{SR-SYS-003}} & {\href{../07\_Программные\_требования/DOC-06.2\_Сессии\_задач\_и\_планировщик.pdf\#swr-ses-001}{SWR-SES-001}}, {\href{../07\_Программные\_требования/DOC-06.2\_Сессии\_задач\_и\_планировщик.pdf\#swr-ses-002}{SWR-SES-002}}, {\href{../07\_Программные\_требования/DOC-06\_Программные\_требования.pdf\#swr-sys-001}{SWR-SYS-001}} \\
{\hyperlink{tc-ses-003}{TC-SES-003}} & {\href{../04\_Системные\_требования/DOC-05\_Системные\_требования.pdf\#sr-ses-002}{SR-SES-002}}, {\href{../04\_Системные\_требования/DOC-05\_Системные\_требования.pdf\#sr-ses-014}{SR-SES-014}}, {\href{../04\_Системные\_требования/DOC-05\_Системные\_требования.pdf\#sr-sys-008}{SR-SYS-008}} & {\href{../07\_Программные\_требования/DOC-06.2\_Сессии\_задач\_и\_планировщик.pdf\#swr-ses-003}{SWR-SES-003}}, {\href{../07\_Программные\_требования/DOC-06.2\_Сессии\_задач\_и\_планировщик.pdf\#swr-ses-022}{SWR-SES-022}}, {\href{../07\_Программные\_требования/DOC-06.2\_Сессии\_задач\_и\_планировщик.pdf\#swr-ses-023}{SWR-SES-023}} \\
{\hyperlink{tc-ses-004}{TC-SES-004}} & {\href{../04\_Системные\_требования/DOC-05\_Системные\_требования.pdf\#sr-ses-001}{SR-SES-001}}, {\href{../04\_Системные\_требования/DOC-05\_Системные\_требования.pdf\#sr-sys-003}{SR-SYS-003}}, {\href{../04\_Системные\_требования/DOC-05\_Системные\_требования.pdf\#sr-sys-005}{SR-SYS-005}} & {\href{../07\_Программные\_требования/DOC-06.2\_Сессии\_задач\_и\_планировщик.pdf\#swr-ses-004}{SWR-SES-004}}, {\href{../07\_Программные\_требования/DOC-06.2\_Сессии\_задач\_и\_планировщик.pdf\#swr-ses-005}{SWR-SES-005}} \\
{\hyperlink{tc-ses-005}{TC-SES-005}} & {\href{../04\_Системные\_требования/DOC-05\_Системные\_требования.pdf\#sr-ses-009}{SR-SES-009}}, {\href{../04\_Системные\_требования/DOC-05\_Системные\_требования.pdf\#sr-sys-008}{SR-SYS-008}}, {\href{../04\_Системные\_требования/DOC-05\_Системные\_требования.pdf\#sr-sys-009}{SR-SYS-009}} & {\href{../07\_Программные\_требования/DOC-06.2\_Сессии\_задач\_и\_планировщик.pdf\#swr-ses-009}{SWR-SES-009}}, {\href{../07\_Программные\_требования/DOC-06.2\_Сессии\_задач\_и\_планировщик.pdf\#swr-ses-014}{SWR-SES-014}}, {\href{../07\_Программные\_требования/DOC-06.2\_Сессии\_задач\_и\_планировщик.pdf\#swr-ses-022}{SWR-SES-022}}, {\href{../07\_Программные\_требования/DOC-06\_Программные\_требования.pdf\#swr-sys-007}{SWR-SYS-007}} \\
{\hyperlink{tc-ses-006}{TC-SES-006}} & {\href{../04\_Системные\_требования/DOC-05\_Системные\_требования.pdf\#sr-ses-004}{SR-SES-004}}, {\href{../04\_Системные\_требования/DOC-05\_Системные\_требования.pdf\#sr-ses-005}{SR-SES-005}}, {\href{../04\_Системные\_требования/DOC-05\_Системные\_требования.pdf\#sr-ses-007}{SR-SES-007}}, {\href{../04\_Системные\_требования/DOC-05\_Системные\_требования.pdf\#sr-sys-002}{SR-SYS-002}}, {\href{../04\_Системные\_требования/DOC-05\_Системные\_требования.pdf\#sr-sys-004}{SR-SYS-004}}, {\href{../04\_Системные\_требования/DOC-05\_Системные\_требования.pdf\#sr-sys-005}{SR-SYS-005}}, {\href{../04\_Системные\_требования/DOC-05\_Системные\_требования.pdf\#sr-sys-009}{SR-SYS-009}} & {\href{../07\_Программные\_требования/DOC-06.2\_Сессии\_задач\_и\_планировщик.pdf\#swr-ses-010}{SWR-SES-010}}, {\href{../07\_Программные\_требования/DOC-06.2\_Сессии\_задач\_и\_планировщик.pdf\#swr-ses-011}{SWR-SES-011}}, {\href{../07\_Программные\_требования/DOC-06.2\_Сессии\_задач\_и\_планировщик.pdf\#swr-ses-012}{SWR-SES-012}}, {\href{../07\_Программные\_требования/DOC-06.2\_Сессии\_задач\_и\_планировщик.pdf\#swr-ses-018}{SWR-SES-018}}, {\href{../07\_Программные\_требования/DOC-06\_Программные\_требования.pdf\#swr-sys-002}{SWR-SYS-002}}, {\href{../07\_Программные\_требования/DOC-06\_Программные\_требования.pdf\#swr-sys-003}{SWR-SYS-003}}, {\href{../07\_Программные\_требования/DOC-06\_Программные\_требования.pdf\#swr-sys-004}{SWR-SYS-004}}, {\href{../07\_Программные\_требования/DOC-06\_Программные\_требования.pdf\#swr-sys-005}{SWR-SYS-005}} \\
{\hyperlink{tc-ses-007}{TC-SES-007}} & {\href{../04\_Системные\_требования/DOC-05\_Системные\_требования.pdf\#sr-ses-008}{SR-SES-008}}, {\href{../04\_Системные\_требования/DOC-05\_Системные\_требования.pdf\#sr-sys-005}{SR-SYS-005}} & {\href{../07\_Программные\_требования/DOC-06.2\_Сессии\_задач\_и\_планировщик.pdf\#swr-ses-008}{SWR-SES-008}}, {\href{../07\_Программные\_требования/DOC-06.2\_Сессии\_задач\_и\_планировщик.pdf\#swr-ses-012}{SWR-SES-012}}, {\href{../07\_Программные\_требования/DOC-06.2\_Сессии\_задач\_и\_планировщик.pdf\#swr-ses-013}{SWR-SES-013}}, {\href{../07\_Программные\_требования/DOC-06.2\_Сессии\_задач\_и\_планировщик.pdf\#swr-ses-021}{SWR-SES-021}} \\
{\hyperlink{tc-ses-008}{TC-SES-008}} & {\href{../04\_Системные\_требования/DOC-05\_Системные\_требования.pdf\#sr-ses-002}{SR-SES-002}}, {\href{../04\_Системные\_требования/DOC-05\_Системные\_требования.pdf\#sr-ses-003}{SR-SES-003}}, {\href{../04\_Системные\_требования/DOC-05\_Системные\_требования.pdf\#sr-ses-006}{SR-SES-006}}, {\href{../04\_Системные\_требования/DOC-05\_Системные\_требования.pdf\#sr-ses-014}{SR-SES-014}}, {\href{../04\_Системные\_требования/DOC-05\_Системные\_требования.pdf\#sr-sys-001}{SR-SYS-001}}, {\href{../04\_Системные\_требования/DOC-05\_Системные\_требования.pdf\#sr-sys-009}{SR-SYS-009}} & {\href{../07\_Программные\_требования/DOC-06.2\_Сессии\_задач\_и\_планировщик.pdf\#swr-ses-001}{SWR-SES-001}}, {\href{../07\_Программные\_требования/DOC-06.2\_Сессии\_задач\_и\_планировщик.pdf\#swr-ses-003}{SWR-SES-003}}, {\href{../07\_Программные\_требования/DOC-06.2\_Сессии\_задач\_и\_планировщик.pdf\#swr-ses-008}{SWR-SES-008}}, {\href{../07\_Программные\_требования/DOC-06.2\_Сессии\_задач\_и\_планировщик.pdf\#swr-ses-017}{SWR-SES-017}}, {\href{../07\_Программные\_требования/DOC-06.2\_Сессии\_задач\_и\_планировщик.pdf\#swr-ses-021}{SWR-SES-021}}, {\href{../07\_Программные\_требования/DOC-06.2\_Сессии\_задач\_и\_планировщик.pdf\#swr-ses-023}{SWR-SES-023}}, {\href{../07\_Программные\_требования/DOC-06\_Программные\_требования.pdf\#swr-sys-001}{SWR-SYS-001}} \\
{\hyperlink{tc-ses-009}{TC-SES-009}} & {\href{../04\_Системные\_требования/DOC-05\_Системные\_требования.pdf\#sr-ses-002}{SR-SES-002}}, {\href{../04\_Системные\_требования/DOC-05\_Системные\_требования.pdf\#sr-ses-006}{SR-SES-006}}, {\href{../04\_Системные\_требования/DOC-05\_Системные\_требования.pdf\#sr-ses-007}{SR-SES-007}}, {\href{../04\_Системные\_требования/DOC-05\_Системные\_требования.pdf\#sr-ses-011}{SR-SES-011}}, {\href{../04\_Системные\_требования/DOC-05\_Системные\_требования.pdf\#sr-sys-006}{SR-SYS-006}} & {\href{../07\_Программные\_требования/DOC-06.2\_Сессии\_задач\_и\_планировщик.pdf\#swr-ses-003}{SWR-SES-003}}, {\href{../07\_Программные\_требования/DOC-06.2\_Сессии\_задач\_и\_планировщик.pdf\#swr-ses-007}{SWR-SES-007}}, {\href{../07\_Программные\_требования/DOC-06.2\_Сессии\_задач\_и\_планировщик.pdf\#swr-ses-015}{SWR-SES-015}}, {\href{../07\_Программные\_требования/DOC-06.2\_Сессии\_задач\_и\_планировщик.pdf\#swr-ses-017}{SWR-SES-017}}, {\href{../07\_Программные\_требования/DOC-06.2\_Сессии\_задач\_и\_планировщик.pdf\#swr-ses-018}{SWR-SES-018}}, {\href{../07\_Программные\_требования/DOC-06\_Программные\_требования.pdf\#swr-sys-006}{SWR-SYS-006}} \\
{\hyperlink{tc-ses-010}{TC-SES-010}} & {\href{../04\_Системные\_требования/DOC-05\_Системные\_требования.pdf\#sr-ses-006}{SR-SES-006}}, {\href{../04\_Системные\_требования/DOC-05\_Системные\_требования.pdf\#sr-ses-011}{SR-SES-011}}, {\href{../04\_Системные\_требования/DOC-05\_Системные\_требования.pdf\#sr-sys-006}{SR-SYS-006}} & {\href{../07\_Программные\_требования/DOC-06.2\_Сессии\_задач\_и\_планировщик.pdf\#swr-ses-007}{SWR-SES-007}}, {\href{../07\_Программные\_требования/DOC-06.2\_Сессии\_задач\_и\_планировщик.pdf\#swr-ses-008}{SWR-SES-008}}, {\href{../07\_Программные\_требования/DOC-06.2\_Сессии\_задач\_и\_планировщик.pdf\#swr-ses-016}{SWR-SES-016}}, {\href{../07\_Программные\_требования/DOC-06.2\_Сессии\_задач\_и\_планировщик.pdf\#swr-ses-017}{SWR-SES-017}}, {\href{../07\_Программные\_требования/DOC-06.2\_Сессии\_задач\_и\_планировщик.pdf\#swr-ses-023}{SWR-SES-023}}, {\href{../07\_Программные\_требования/DOC-06\_Программные\_требования.pdf\#swr-sys-006}{SWR-SYS-006}} \\
{\hyperlink{tc-ses-011}{TC-SES-011}} & {\href{../04\_Системные\_требования/DOC-05\_Системные\_требования.pdf\#sr-ses-010}{SR-SES-010}}, {\href{../04\_Системные\_требования/DOC-05\_Системные\_требования.pdf\#sr-ses-012}{SR-SES-012}} & {\href{../07\_Программные\_требования/DOC-06.2\_Сессии\_задач\_и\_планировщик.pdf\#swr-ses-014}{SWR-SES-014}}, {\href{../07\_Программные\_требования/DOC-06.2\_Сессии\_задач\_и\_планировщик.pdf\#swr-ses-020}{SWR-SES-020}} \\
{\hyperlink{tc-ses-012}{TC-SES-012}} & {\href{../04\_Системные\_требования/DOC-05\_Системные\_требования.pdf\#sr-ses-009}{SR-SES-009}}, {\href{../04\_Системные\_требования/DOC-05\_Системные\_требования.pdf\#sr-ses-010}{SR-SES-010}}, {\href{../04\_Системные\_требования/DOC-05\_Системные\_требования.pdf\#sr-ses-011}{SR-SES-011}}, {\href{../04\_Системные\_требования/DOC-05\_Системные\_требования.pdf\#sr-sys-009}{SR-SYS-009}} & {\href{../07\_Программные\_требования/DOC-06.2\_Сессии\_задач\_и\_планировщик.pdf\#swr-ses-009}{SWR-SES-009}}, {\href{../07\_Программные\_требования/DOC-06.2\_Сессии\_задач\_и\_планировщик.pdf\#swr-ses-014}{SWR-SES-014}}, {\href{../07\_Программные\_требования/DOC-06.2\_Сессии\_задач\_и\_планировщик.pdf\#swr-ses-015}{SWR-SES-015}}, {\href{../07\_Программные\_требования/DOC-06.2\_Сессии\_задач\_и\_планировщик.pdf\#swr-ses-017}{SWR-SES-017}}, {\href{../07\_Программные\_требования/DOC-06\_Программные\_требования.pdf\#swr-sys-004}{SWR-SYS-004}}, {\href{../07\_Программные\_требования/DOC-06\_Программные\_требования.pdf\#swr-sys-007}{SWR-SYS-007}} \\
{\hyperlink{tc-ses-013}{TC-SES-013}} & {\href{../04\_Системные\_требования/DOC-05\_Системные\_требования.pdf\#sr-ses-008}{SR-SES-008}}, {\href{../04\_Системные\_требования/DOC-05\_Системные\_требования.pdf\#sr-ses-012}{SR-SES-012}}, {\href{../04\_Системные\_требования/DOC-05\_Системные\_требования.pdf\#sr-ses-013}{SR-SES-013}}, {\href{../04\_Системные\_требования/DOC-05\_Системные\_требования.pdf\#sr-sys-007}{SR-SYS-007}} & {\href{../07\_Программные\_требования/DOC-06.2\_Сессии\_задач\_и\_планировщик.pdf\#swr-ses-019}{SWR-SES-019}}, {\href{../07\_Программные\_требования/DOC-06.2\_Сессии\_задач\_и\_планировщик.pdf\#swr-ses-020}{SWR-SES-020}}, {\href{../07\_Программные\_требования/DOC-06.2\_Сессии\_задач\_и\_планировщик.pdf\#swr-ses-021}{SWR-SES-021}} \\
{\hyperlink{tc-ses-014}{TC-SES-014}} & {\href{../04\_Системные\_требования/DOC-05\_Системные\_требования.pdf\#sr-ses-005}{SR-SES-005}}, {\href{../04\_Системные\_требования/DOC-05\_Системные\_требования.pdf\#sr-ses-008}{SR-SES-008}}, {\href{../04\_Системные\_требования/DOC-05\_Системные\_требования.pdf\#sr-ses-014}{SR-SES-014}}, {\href{../04\_Системные\_требования/DOC-05\_Системные\_требования.pdf\#sr-sys-002}{SR-SYS-002}}, {\href{../04\_Системные\_требования/DOC-05\_Системные\_требования.pdf\#sr-sys-007}{SR-SYS-007}}, {\href{../04\_Системные\_требования/DOC-05\_Системные\_требования.pdf\#sr-sys-008}{SR-SYS-008}} & {\href{../07\_Программные\_требования/DOC-06.2\_Сессии\_задач\_и\_планировщик.pdf\#swr-ses-011}{SWR-SES-011}}, {\href{../07\_Программные\_требования/DOC-06.2\_Сессии\_задач\_и\_планировщик.pdf\#swr-ses-013}{SWR-SES-013}}, {\href{../07\_Программные\_требования/DOC-06.2\_Сессии\_задач\_и\_планировщик.pdf\#swr-ses-019}{SWR-SES-019}}, {\href{../07\_Программные\_требования/DOC-06.2\_Сессии\_задач\_и\_планировщик.pdf\#swr-ses-021}{SWR-SES-021}}, {\href{../07\_Программные\_требования/DOC-06.2\_Сессии\_задач\_и\_планировщик.pdf\#swr-ses-023}{SWR-SES-023}}, {\href{../07\_Программные\_требования/DOC-06\_Программные\_требования.pdf\#swr-sys-008}{SWR-SYS-008}} \\
{\hyperlink{tc-ses-015}{TC-SES-015}} & {\href{../04\_Системные\_требования/DOC-05\_Системные\_требования.pdf\#sr-ses-005}{SR-SES-005}}, {\href{../04\_Системные\_требования/DOC-05\_Системные\_требования.pdf\#sr-ses-010}{SR-SES-010}}, {\href{../04\_Системные\_требования/DOC-05\_Системные\_требования.pdf\#sr-sys-005}{SR-SYS-005}}, {\href{../04\_Системные\_требования/DOC-05\_Системные\_требования.pdf\#sr-sys-009}{SR-SYS-009}} & {\href{../07\_Программные\_требования/DOC-06.2\_Сессии\_задач\_и\_планировщик.pdf\#swr-ses-011}{SWR-SES-011}}, {\href{../07\_Программные\_требования/DOC-06.2\_Сессии\_задач\_и\_планировщик.pdf\#swr-ses-014}{SWR-SES-014}}, {\href{../07\_Программные\_требования/DOC-06.2\_Сессии\_задач\_и\_планировщик.pdf\#swr-ses-021}{SWR-SES-021}}, {\href{../07\_Программные\_требования/DOC-06\_Программные\_требования.pdf\#swr-sys-005}{SWR-SYS-005}}, {\href{../07\_Программные\_требования/DOC-06\_Программные\_требования.pdf\#swr-sys-007}{SWR-SYS-007}} \\
{\hyperlink{tc-ses-016}{TC-SES-016}} & {\href{../04\_Системные\_требования/DOC-05\_Системные\_требования.pdf\#sr-ses-007}{SR-SES-007}}, {\href{../04\_Системные\_требования/DOC-05\_Системные\_требования.pdf\#sr-ses-008}{SR-SES-008}}, {\href{../04\_Системные\_требования/DOC-05\_Системные\_требования.pdf\#sr-ses-011}{SR-SES-011}}, {\href{../04\_Системные\_требования/DOC-05\_Системные\_требования.pdf\#sr-sys-005}{SR-SYS-005}} & {\href{../07\_Программные\_требования/DOC-06.2\_Сессии\_задач\_и\_планировщик.pdf\#swr-ses-011}{SWR-SES-011}}, {\href{../07\_Программные\_требования/DOC-06.2\_Сессии\_задач\_и\_планировщик.pdf\#swr-ses-012}{SWR-SES-012}}, {\href{../07\_Программные\_требования/DOC-06.2\_Сессии\_задач\_и\_планировщик.pdf\#swr-ses-015}{SWR-SES-015}}, {\href{../07\_Программные\_требования/DOC-06.2\_Сессии\_задач\_и\_планировщик.pdf\#swr-ses-018}{SWR-SES-018}}, {\href{../07\_Программные\_требования/DOC-06.2\_Сессии\_задач\_и\_планировщик.pdf\#swr-ses-021}{SWR-SES-021}}, {\href{../07\_Программные\_требования/DOC-06\_Программные\_требования.pdf\#swr-sys-005}{SWR-SYS-005}} \\
{\hyperlink{tc-ses-017}{TC-SES-017}} & {\href{../04\_Системные\_требования/DOC-05\_Системные\_требования.pdf\#sr-cfg-005}{SR-CFG-005}}, {\href{../04\_Системные\_требования/DOC-05\_Системные\_требования.pdf\#sr-ses-006}{SR-SES-006}}, {\href{../04\_Системные\_требования/DOC-05\_Системные\_требования.pdf\#sr-ses-011}{SR-SES-011}}, {\href{../04\_Системные\_требования/DOC-05\_Системные\_требования.pdf\#sr-ses-012}{SR-SES-012}}, {\href{../04\_Системные\_требования/DOC-05\_Системные\_требования.pdf\#sr-sys-006}{SR-SYS-006}} & {\href{../07\_Программные\_требования/DOC-06\_Программные\_требования.pdf\#swr-cfg-005}{SWR-CFG-005}}, {\href{../07\_Программные\_требования/DOC-06.2\_Сессии\_задач\_и\_планировщик.pdf\#swr-ses-006}{SWR-SES-006}}, {\href{../07\_Программные\_требования/DOC-06.2\_Сессии\_задач\_и\_планировщик.pdf\#swr-ses-007}{SWR-SES-007}}, {\href{../07\_Программные\_требования/DOC-06.2\_Сессии\_задач\_и\_планировщик.pdf\#swr-ses-015}{SWR-SES-015}}, {\href{../07\_Программные\_требования/DOC-06.2\_Сессии\_задач\_и\_планировщик.pdf\#swr-ses-017}{SWR-SES-017}}, {\href{../07\_Программные\_требования/DOC-06.2\_Сессии\_задач\_и\_планировщик.pdf\#swr-ses-020}{SWR-SES-020}}, {\href{../07\_Программные\_требования/DOC-06\_Программные\_требования.pdf\#swr-sys-006}{SWR-SYS-006}} \\
\end{longtable}
}


\end{document}
